% !TeX program = xelatex
\documentclass[12pt]{article}
\usepackage{makecell}
\usepackage{hyphenat}
\usepackage{fontspec}
% Load the TeX Live font files directly so the document also builds on
% machines where the system font index does not expose "TeX Gyre Termes".
\setmainfont{texgyretermes-regular.otf}[
  BoldFont=texgyretermes-bold.otf,
  ItalicFont=texgyretermes-italic.otf,
  BoldItalicFont=texgyretermes-bolditalic.otf
]
\usepackage{polyglossia}
\setdefaultlanguage{vietnamese}\setotherlanguage{english}
\usepackage[table,dvipsnames]{xcolor} % Required for row coloring
\usepackage{booktabs} % Required for nicer horizontal lines
\usepackage{ragged2e}
\usepackage{caption}
\usepackage{subcaption}
\usepackage{float}
\usepackage{amsmath,amssymb,amsfonts,amsthm}
\usepackage{graphicx}
\usepackage{adjustbox}
\usepackage{multicol}

\newtheorem{theorem}{Định lý}
\newtheorem{lemma}{Bổ đề}
\newtheorem{proposition}{Mệnh đề}
\newtheorem{corollary}{Hệ quả}
\newtheorem{definition}{Định nghĩa}
\newtheorem{remark}{Nhận xét}
\usepackage[inline,shortlabels]{enumitem}
\usepackage{setspace}
\usepackage{longtable}
\usepackage{bm}
\usepackage{titling}
\usepackage{tabularx}
% fontenc is not needed with XeLaTeX
\usepackage{mathrsfs}
% \usepackage[draft]{hyperref}
\usepackage{hyperref}
\usepackage{xparse}
\usepackage{fontawesome5}
\usepackage[most]{tcolorbox}
\usepackage{tikz}
\usepackage[framemethod=TikZ]{mdframed}
\usetikzlibrary{matrix, fit, shadows, positioning, arrows.meta, calc, shapes.geometric, backgrounds, shapes.symbols}
\usepackage{paracol}
\usepackage{siunitx}
\usepackage{multirow}
\usepackage{forest}

\usepackage[toc,page]{appendix}
\renewcommand{\appendixtocname}{Phụ lục}
\renewcommand{\appendixpagename}{Phụ lục}

\usepackage{xurl}
\setlength{\headheight}{14.5pt}
\setlength{\topmargin}{-.5in} \setlength{\textheight}{9.25in}
\setlength{\oddsidemargin}{-0.1in} \setlength{\textwidth}{6.8in}

\usepackage{tocloft}
\renewcommand{\cfttoctitlefont}{\Huge\bfseries}
\setlength{\cftaftertoctitleskip}{2em}
\setlength{\cftbeforesubsecskip}{7pt} % giãn cách subitem mục lục
\renewcommand{\cftsecleader}{\cftdotfill{\cftdotsep}}
\setlength{\cftsecnumwidth}{3.2em}
\renewcommand{\cftsecaftersnum}{\quad}
\renewcommand{\cftsubsecleader}{\cftdotfill{\cftdotsep}}
\setlength{\cftsubsecnumwidth}{4em}
\renewcommand{\cftsubsecaftersnum}{\quad}
\renewcommand{\cftsubsubsecleader}{\cftdotfill{\cftdotsep}}
\setlength   {\cftsubsubsecnumwidth}{5em}
\renewcommand{\cftsubsubsecaftersnum}{\quad}

\newcounter{mycounter}
\newcommand\showmycounter{\stepcounter{mycounter}\themycounter}
\usepackage{lipsum}
\newcommand\showlips{\stepcounter{mycounter}\lipsum[\value{mycounter}]}

%--------------------------------------------------------------------
\usepackage{framed}
\usepackage{fancyhdr}
\usepackage{upquote}
\usepackage{listings}


% Color --------------------------------------------------------------------
\usepackage{epigraph}
\usepackage{titlesec}

\definecolor{codebackground}{RGB}{250, 250, 250}
\definecolor{codecomment}{RGB}{34, 139, 34}
\definecolor{codekeyword}{RGB}{0, 0, 255}
\definecolor{codestring}{RGB}{255, 0, 0}
\definecolor{codeidentifier}{RGB}{0, 0, 0}
\definecolor{codelinenumbers}{RGB}{90, 90, 90}

\definecolor{codegreen}{rgb}{0,0.6,0}
\definecolor{codegray}{rgb}{0.5,0.5,0.5}
\definecolor{codepurple}{rgb}{0.58,0,0.82}
\definecolor{backcolour}{rgb}{0.95,0.95,0.92}

\definecolor{aivnGreen}{RGB}{52,79,30}
\definecolor{aivnCodeOutputColor}{RGB}{11,83,148}
\definecolor{captiongray}{gray}{0.15}
\definecolor{captionbg}{RGB}{250,250,250}

\definecolor{foldercolor}{RGB}{0, 102, 204}
\definecolor{filecolor}{RGB}{104, 104, 104}
\definecolor{scriptcolor}{RGB}{180, 50, 50}
\definecolor{pythoncolor}{RGB}{0, 128, 128}

% --- Extra Brand colors for the table ---
\definecolor{brandIndigo}{HTML}{333D91}
\definecolor{rowLavender}{HTML}{EEECFA}
\setlength{\parskip}{0.3em}

% Table
\newtcolorbox{captionbelowboxaivn}{
    colback=aivnGreen,
    colframe=aivnGreen,
    boxrule=0pt,
    arc=4pt,
    left=4pt,
    right=4pt,
    top=2pt,
    bottom=2pt,
    width=\linewidth,
    boxsep=4pt,
    halign=center,
    % fontupper=\bfseries\color{white}
    fontupper=\color{white}
}

% Figure
\newtcolorbox{figurecardcaptionaivn}{
  colback=aivnGreen,        % Màu nền (background color)
  colframe=aivnGreen,       % Màu viền box (border color)
  boxrule=0pt,              % Độ dày viền box (0pt = không viền)
  arc=2pt,                  % Độ bo góc (2pt là bo nhẹ, tăng để bo tròn hơn)
  top=3pt,                  % Padding phía trên nội dung
  bottom=3pt,               % Padding phía dưới nội dung
  left=10pt,                % Padding bên trái
  right=10pt,               % Padding bên phải
  width=\linewidth,         % Chiều rộng bằng toàn bộ dòng
  fontupper=\small\fontsize{11pt}{13pt}\selectfont\color{white}
}

% Tạo counter riêng
\newcounter{hinh}
% Định nghĩa macro caption Hình
\newcommand{\capfig}[1]{%
  \refstepcounter{hinh}%
  % \textbf{Hình~\thehinh:}~#1
    Hình~\thehinh:~#1
}

\newtcolorbox{figurecaptionmodern}{
  colback=captionbg,
  colframe=aivnGreen,
  boxrule=1.5pt,
  bottomrule=2pt,
  toptitle=0pt,
  bottomtitle=0pt,
  left=6pt,
  right=6pt,
  top=4pt,
  bottom=2pt,
  width=\linewidth,
  fontupper=\small\fontsize{11pt}{13pt}\selectfont\color{black},
  halign=center
}

% Code Block -----------------------------------------------------------------
\newtcolorbox{aivncodebox}[1][]{%
  enhanced,
  colback=gray!10,
  colframe=aivnGreen,
  coltitle=white,
  colbacktitle=aivnGreen,
  title=#1,
  fonttitle=\bfseries\sffamily,
  listing only,
  breakable,
  listing engine=listings,
  listing options={
    language=Python,
    basicstyle=\ttfamily\footnotesize,
    commentstyle=\color{ForestGreen},
    keywordstyle=\color{blue},
    stringstyle=\color{red!70!black},
    numbers=left, upquote=true,
    numberstyle=\tiny\color{gray},
    numbersep=5pt,
    breaklines=true,
    showstringspaces=false,
    xleftmargin=1.5em,
    framexleftmargin=1.5em
  }
}

% Code Output Block -----------------------------------------------------------------
\newtcolorbox{aivnCodeOutput}[1][]{%
  enhanced,
  colback=gray!10,
  colframe=aivnCodeOutputColor,
  coltitle=white,
  colbacktitle=aivnCodeOutputColor,
  title=#1,
  fonttitle=\bfseries\sffamily,
  listing only,
  breakable,
  listing engine=listings,
  listing options={
    language=Python,
    basicstyle=\ttfamily\footnotesize,
    commentstyle=\color{ForestGreen},
    keywordstyle=\color{blue},
    stringstyle=\color{red!70!black},
    numbers=left,
    numberstyle=\tiny\color{gray},
    numbersep=5pt,
    breaklines=true,
    showstringspaces=false,
    xleftmargin=1.5em,
    framexleftmargin=1.5em
  }
}

% Last argument allows more tcolorbox options to be added

% --- Guideline/Tip box ---
\newtcolorbox{guidelinebox}[2][]{
    enhanced, colback=yellow!10!white, colframe=orange!90!black, boxrule=1pt, arc=4pt,
    fonttitle=\bfseries, attach boxed title to top left={yshift=-2mm, xshift=3mm},
    colbacktitle=orange!90!black, coltitle=white,
    title=\faLightbulb\hspace{1mm} #2, #1
}

\newtcolorbox{notebox}[1][]{
  enhanced,  colback=yellow!10, colframe=orange!80, boxrule=1pt,
  arc=4pt, fontupper=\small, attach boxed title to top left={yshift=-2mm, xshift=3mm},
  colbacktitle=orange!80,  coltitle=white, fonttitle=\bfseries,
  title=\faInfoCircle~\textbf{Lưu ý},  #1
}

\newtcolorbox{infobox}[1]{%
    enhanced,
    breakable,
    colback=green!3,
    colframe=green!3,
    arc=5pt,
    borderline west={3pt}{0pt}{aivnGreen},
    attach boxed title to top left={
        yshift*=-\tcboxedtitleheight/2,
        xshift=-0.1em
    },
    colbacktitle=aivnGreen, %
    coltitle=white,
    fonttitle=\bfseries,
    boxed title style={
        arc=4pt,
        top=2pt,
        bottom=2pt,
        left=8pt,
        right=8pt,
    },
    title={\faInfoCircle\quad #1}
}
\newtcolorbox{examplebox}[1]{%
    enhanced,
    breakable,
    colback=brandIndigo!5,
    colframe=brandIndigo!5,
    arc=5pt,
    borderline west={3pt}{0pt}{brandIndigo},
    attach boxed title to top left={
        yshift*=-\tcboxedtitleheight/2,
        xshift=-0.1em
    },
    colbacktitle=brandIndigo,
    coltitle=white,
    fonttitle=\bfseries,
    boxed title style={
        arc=4pt,
        top=2pt,
        bottom=2pt,
        left=8pt,
        right=8pt,
    },
    title={\faLightbulb\quad #1}
}
% Prompt Block ---------------------------------------------------------------
\definecolor{promptBlue}{HTML}{2F5D8C}
\definecolor{promptBg}{HTML}{F6F9FC}
\definecolor{promptBorder}{HTML}{B7C9DC}

% Dùng cho prompt đã viết theo cú pháp LaTeX bình thường
\newtcolorbox{promptbox}[1][Prompt]{%
    enhanced,
    breakable,
    colback=promptBg,
    colframe=promptBorder,
    coltitle=white,
    colbacktitle=promptBlue,
    fonttitle=\bfseries\sffamily,
    title={#1},
    boxrule=0.8pt,
    arc=4pt,
    left=8pt,
    right=8pt,
    top=2pt,
    bottom=1pt,
    before skip=8pt,
    after skip=8pt,
    borderline west={3pt}{0pt}{promptBlue},
    fontupper=\small\sffamily,
}

% Dùng cho prompt raw: giữ nguyên xuống dòng, không highlight keyword
\newtcblisting{promptboxraw}[1][Prompt]{%
    enhanced,
    breakable,
    listing only,
    listing engine=listings,
    colback=promptBg,
    colframe=promptBorder,
    coltitle=white,
    colbacktitle=promptBlue,
    fonttitle=\bfseries\sffamily,
    title={#1},
    boxrule=0.8pt,
    arc=4pt,
    left=8pt,
    right=8pt,
    top=-4pt,
    bottom=-4pt,
    before skip=4pt,
    after skip=4pt,
    borderline west={3pt}{0pt}{promptBlue},
    listing options={
        language={},
        basicstyle=\small\ttfamily,
        columns=fullflexible,
        keepspaces=true,
        breaklines=true,
        showstringspaces=false,
        numbers=none,
        frame=none,
        keywordstyle=\color{black},
        commentstyle=\color{black},
        stringstyle=\color{black},
        identifierstyle=\color{black}
    }
}


%--------------------------------------------------------------------
\lstdefinestyle{mystyle}{
    backgroundcolor=\color{codebackground},
    commentstyle=\color{codecomment},
    keywordstyle=\color{codekeyword},
    stringstyle=\color{codestring},
    identifierstyle=\color{codeidentifier},
    basicstyle=\ttfamily\footnotesize,
    breakatwhitespace=false,
    breaklines=true,
    captionpos=b,
    keepspaces=true,
    numbers=left,
    numberstyle=\tiny\color{codelinenumbers},
    numbersep=5pt,
    showspaces=false,
    showstringspaces=false,
    showtabs=false,
    tabsize=2
}

\lstset{style=mystyle}

\hypersetup{
    colorlinks=true,
    linkcolor=black,
    filecolor=magenta,
    citecolor=green,
    urlcolor=blue,
    pdftitle={Nguyên nhân và cơ chế của Vanishing Gradient},
    pdfauthor={Tâm, Trang, Trí, Khanh, Khoa, Long},
    pdflang={vi},
    pdfpagemode=FullScreen,
}

%--------------------------------------------------------------------
\newcounter{choice}
\renewcommand\thechoice{($\alph{choice}$)}
\newcommand\choicelabel{\thechoice}

%--------------------------------------------------------------------
\newenvironment{choices}%
  {\list{\choicelabel}%
     {\usecounter{choice}\def\makelabel##1{\hss\llap{##1}}%
       \settowidth{\leftmargin}{W.\hskip\labelsep\hskip 0 em}%
       \def\choice{%
         \item
       } % choice
       \labelwidth\leftmargin\advance\labelwidth-\labelsep
       \topsep=0pt
       \partopsep=0pt
     }%
  }%
  {\endlist}

%--------------------------------------------------------------------
\def\CChoice{%
      \choice
      \addanswer{\theenumi}{\thechoice}%
    }
    \let\CChoice\CChoice
%    \par % Uncomment this to have choices always start a new line
   % \let\par\@empty
    % If we're continuing the textbf containing the question,
    % then leave a bit of space before the first choice:
    \ifvmode\else\enskip\fi
    \ignorespaces
  %
  {}
\makeatother
\newbox\allanswers
\setbox\allanswers=\hbox{}
\newcommand{\addanswer}[2]{%
  \global\setbox\allanswers=\hbox{\unhbox\allanswers \quad #1.~#2}%
}
\newcommand{\showanswers}{%
  \vfill
  \begin{center}
    Answers
  \end{center}
  \unhbox\allanswers
}

\makeatletter
\newenvironment{multichoices}[1][1]{%
\begin{multicols}{#1}}{%
\end{multicols}}
\makeatother

\pagestyle{fancy}

\DeclareMathOperator*{\argmax}{arg\,max}
\DeclareMathOperator*{\argmin}{arg\,min}
\DeclareMathOperator{\E}{\mathbb{E}}

% Outline --------------------------------------------------------------------
% Section numbering ----------------------------------------------------------
\renewcommand{\thesection}{\Roman{section}.}
\renewcommand{\thesubsection}{\thesection\arabic{subsection}.}
\renewcommand{\thesubsubsection}{\thesubsection\arabic{subsubsection}.}

% Paper-like section title sizes --------------------------------------------
% Title paper: 17pt
% Section should be clearly smaller than title but still prominent.
\titleformat{\section}
  {\normalfont\fontsize{16pt}{20pt}\selectfont\bfseries}
  {\thesection}
  {0.85em}
  {}

\titleformat{\subsection}
  {\normalfont\fontsize{14pt}{16pt}\selectfont\bfseries}
  {\thesubsection}
  {0.75em}
  {}

\titleformat{\subsubsection}
  {\normalfont\fontsize{12pt}{15pt}\selectfont\bfseries}
  {\thesubsubsection}
  {0.65em}
  {}

% Paper-like compact spacing -------------------------------------------------
\titlespacing*{\section}
  {0pt}
  {1.25em plus 0.2em minus 0.1em}
  {0.65em plus 0.1em minus 0.1em}

\titlespacing*{\subsection}
  {0pt}
  {1.0em plus 0.2em minus 0.1em}
  {0.45em plus 0.1em minus 0.1em}

\titlespacing*{\subsubsection}
  {0pt}
  {0.85em plus 0.15em minus 0.1em}
  {0.35em plus 0.1em minus 0.1em}


% Logo + Title ---------------------------------------------------------------

\setlength{\droptitle}{-7em}

% Title horizontal rule
\newcommand{\TitleRule}{%
  \rule{0.90\textwidth}{0.6pt}
}

% Chiều cao vùng title giữa 2 rule
\newlength{\TitleBlockHeight}
\setlength{\TitleBlockHeight}{2.2cm}

\pretitle{%


  % Logo row
  \begin{tabular}{@{}c@{\hspace{1.0em}}c@{\hspace{1.0em}}c@{}}
    \raisebox{-0.35\height}{\includegraphics[height=2.0cm]{assets/aio2026_logo.pdf}}
  \end{tabular}

  %\vspace{0.2em}

\begin{center}
  \TitleRule
  \par\nointerlineskip
}

\title{%
  \parbox[c][\TitleBlockHeight][c]{0.90\textwidth}{%
    \centering
    {\fontsize{20pt}{24pt}\selectfont\bfseries Nguyên nhân và cơ chế của Vanishing Gradient}
  }%
}

\posttitle{%
  \par\nointerlineskip
  \TitleRule
  \end{center}
}

\preauthor{%
  \begin{center}
  \fontsize{11pt}{13pt}\selectfont
}

\author{%
  \textbf{Tâm}\hspace{1.6em}
  \textbf{Trang}\hspace{1.6em}
  \textbf{Trí}\hspace{1.6em}
  \textbf{Khanh}\hspace{1.6em}
  \textbf{Khoa}\hspace{1.6em}
  \textbf{Long}\\
  \vspace{1em}
  \textbf{Proofreading: Trang, Trí}
}

\postauthor{%
  \par
  \end{center}
}

\date{}
% Header --------------------------------------------------------------------
\pagestyle{fancy}
\fancyhf{}
\lhead{\href{https://www.facebook.com/aivietnam.edu.vn}{\textcolor{aivnGreen}{\bfseries AI VIET NAM (AIO2026)}}}
\rhead{\href{https://aivietnam.edu.vn/}{\textcolor{aivnGreen}{\bfseries aivietnam.edu.vn}}}
\fancyfoot[r]{\textcolor{aivnGreen}{\bfseries Trang \thepage}}
\fancyfoot[l]{%
  \href{https://zalo.me/84911118758}{%
    \textcolor{aivnGreen}{\bfseries SĐT/Zalo: 0911118758 (Dr. Hồng Phúc)}%
  }%
}
\renewcommand{\headrulewidth}{1.0pt}
\renewcommand{\headrule}{{\color{aivnGreen}\hrule width\headwidth height\headrulewidth \vskip-\headrulewidth}}

\renewcommand{\footrulewidth}{1.0pt}
\renewcommand{\footrule}{{\color{aivnGreen}\hrule width\headwidth height\headrulewidth \vskip-\headrulewidth}}


\fancypagestyle{firstpage}{
    \fancyhf{}
    \fancyfoot[r]{\textcolor{aivnGreen}{\bfseries Trang \thepage}}
    \fancyfoot[l]{%
      \href{https://zalo.me/84911118758}{%
        \textcolor{aivnGreen}{\bfseries SĐT/Zalo: 0911118758 (Dr. Hồng Phúc)}%
      }%
    }
    \renewcommand{\headrulewidth}{0pt}
    \renewcommand{\footrulewidth}{1.0pt}
    \renewcommand{\footrule}{%
        {\color{aivnGreen}\hrule width\headwidth height\footrulewidth}
    }
}


\begin{document}
\maketitle
\thispagestyle{firstpage}
\vspace{-3em}

\tableofcontents
\clearpage

\section{Giới thiệu về Vanishing Gradient}

Trong deep learning nói riêng và học máy nói chung, sự biến mất đạo hàm
(Vanishing Gradient) là một trong những hiện tượng gây ra nhiều khó khăn trong
quá trình huấn luyện các mạng nơ-ron sâu. Hiện tượng này thường xảy ra khi
gradient phải lan truyền ngược qua quá nhiều lớp, hoặc qua một chuỗi tuần tự
dài trong mạng hồi quy (Recurrent Neural Network). Khi gradient truyền đến các
lớp đầu tiên, giá trị cập nhật tham số có thể trở nên cực kỳ nhỏ, khiến các lớp
này gần như không được học trong suốt quá trình huấn luyện. Vì vậy, mô hình khó
nắm bắt các biểu diễn phức tạp của dữ liệu, từ đó bị giới hạn về độ sâu và hiệu
năng.

\begin{figure}[H]
    \centering
    \includegraphics[width=0.80\linewidth]{assets/water_relay_vanishing_gradient.jpg}
    \caption[Vanishing gradient qua trò chơi truyền nước]{Ta có thể tưởng tượng vanishing gradient giống như trò chơi truyền nước mà ta thường chơi ngày bé.
Khi hàng người truyền nước càng dài, hay khoảng cách giữa các người chơi xa,
hoặc người chơi dùng tay để hứng nước thì lượng nước đến được tới thau đựng
nước sẽ ngày càng ít hoặc có thể là không có gì.}
    \label{fig:water-relay-vanishing-gradient}
    {\footnotesize Nguồn ảnh: Fiedler (2024)~\cite{fiedler2024}.\par}
\end{figure}

\section{Kiến trúc nền tảng của mạng Neural Network}
\textit{Người phụ trách: Trang, Trí}

Mạng neural nhân tạo là một hàm tham số hóa $f_{\theta}:\mathbb{R}^{n}\to\mathbb{R}^{K}$ được xây dựng bằng cách xếp chồng nhiều phép biến đổi đơn giản. Mỗi phép biến đổi gồm một bước \emph{tuyến tính} và một bước \emph{phi tuyến}; chính sự kết hợp này cho phép mạng xấp xỉ những quan hệ rất phức tạp giữa đầu vào và đầu ra. Phần này trình bày các khối cơ bản, cách chúng được ghép lại và cách mạng học từ dữ liệu, làm nền tảng để giải thích hiện tượng vanishing gradient ở phần sau.

\subsection{Neuron, trọng số và bias}
Neuron là đơn vị tính toán nhỏ nhất của mạng. Với vector đầu vào $\mathbf{x}\in\mathbb{R}^{n}$, một lớp gồm $m$ neuron trước hết thực hiện phép biến đổi tuyến tính
\begin{equation}
    \mathbf{z}=\mathbf{W}\mathbf{x}+\mathbf{b},
    \qquad \mathbf{W}\in\mathbb{R}^{m\times n},\quad \mathbf{b}\in\mathbb{R}^{m},
    \label{eq:linear}
\end{equation}
rồi đi qua hàm kích hoạt phi tuyến $\sigma$ áp dụng theo từng phần tử để cho đầu ra $\mathbf{a}=\sigma(\mathbf{z})$. Với neuron thứ $j$, ta có $z_j=\sum_{i=1}^{n}W_{ji}x_i+b_j$, tức là tổng có trọng số của các đặc trưng đầu vào cộng thêm một hằng số.

\begin{itemize}
    \item \textbf{Trọng số $\mathbf{W}$} cho biết mức độ và chiều ảnh hưởng của từng đặc trưng lên neuron: trọng số dương khuếch đại, trọng số âm ức chế, trọng số gần $0$ khiến đặc trưng gần như bị bỏ qua. Đây là nhóm tham số chính được điều chỉnh trong quá trình huấn luyện.
    \item \textbf{Bias $\mathbf{b}$} dịch chuyển ngưỡng kích hoạt của neuron, giúp mô hình biểu diễn được các hàm không đi qua gốc tọa độ. Về mặt hình học, $\mathbf{W}$ quyết định hướng của siêu phẳng phân chia, còn $\mathbf{b}$ quyết định vị trí của nó.
    \item \textbf{Hàm kích hoạt $\sigma$} đưa tính phi tuyến vào mạng. Nếu bỏ $\sigma$, tích của nhiều ma trận vẫn là một ma trận, nên mạng sâu chỉ tương đương một phép biến đổi tuyến tính duy nhất.
\end{itemize}

Ba hàm kích hoạt thường gặp, cùng đạo hàm của chúng, được tóm tắt trong Bảng~\ref{tab:activation}. Cột cuối sẽ đặc biệt quan trọng ở Mục~\ref{sec:vanishing}.

\begin{table}[htbp]
    \centering
    \caption{Các hàm kích hoạt phổ biến và độ lớn đạo hàm.}
    \label{tab:activation}
    \begin{tabular}{@{}llcc@{}}
        \toprule
        Hàm & Công thức & Miền giá trị & $\max\sigma'(z)$ \\
        \midrule
        Sigmoid & $\dfrac{1}{1+e^{-z}}$ & $(0,1)$ & $0{,}25$ \\[6pt]
        Tanh & $\dfrac{e^{z}-e^{-z}}{e^{z}+e^{-z}}$ & $(-1,1)$ & $1$ \\[6pt]
        ReLU & $\max(0,z)$ & $[0,\infty)$ & $1$ \\
        \bottomrule
    \end{tabular}
\end{table}

\subsection{Các lớp của mạng neural}
Mạng truyền thẳng (feedforward) được tổ chức thành ba loại lớp:
\begin{itemize}
    \item \textbf{Lớp đầu vào (input layer)} nhận dữ liệu thô $\mathbf{x}$ như điểm ảnh, đặc trưng dạng bảng hay embedding của từ. Lớp này không tính toán mà chỉ đưa dữ liệu vào mạng.
    \item \textbf{Các lớp ẩn (hidden layers)} biến đổi dữ liệu qua nhiều bước liên tiếp. Thông thường, các lớp đầu học đặc trưng đơn giản (cạnh, tần suất từ), còn các lớp sau kết hợp chúng thành biểu diễn trừu tượng hơn (hình dạng, ngữ nghĩa).
    \item \textbf{Lớp đầu ra (output layer)} tạo dự đoán $\hat{\mathbf{y}}$. Hàm kích hoạt của lớp này phụ thuộc bài toán: tuyến tính cho hồi quy, sigmoid cho phân loại nhị phân, softmax cho phân loại nhiều lớp.
\end{itemize}

Thông tin đi từ lớp này sang lớp kế tiếp. Với mạng gồm $L$ lớp (không tính lớp đầu vào) và quy ước $\mathbf{a}^{(0)}=\mathbf{x}$, ta có
\begin{equation}
    \mathbf{z}^{(l)}=\mathbf{W}^{(l)}\mathbf{a}^{(l-1)}+\mathbf{b}^{(l)},
    \qquad
    \mathbf{a}^{(l)}=\sigma^{(l)}\!\bigl(\mathbf{z}^{(l)}\bigr),
    \qquad l=1,\dots,L.
    \label{eq:layer}
\end{equation}
Tập tham số của toàn mạng là $\theta=\{\mathbf{W}^{(l)},\mathbf{b}^{(l)}\}_{l=1}^{L}$.

\subsection{Forward pass và hàm mất mát}
\textbf{Forward pass} là việc áp dụng công thức~\eqref{eq:layer} lần lượt từ lớp $1$ đến lớp $L$ để thu được dự đoán $\hat{\mathbf{y}}=\mathbf{a}^{(L)}$. Trong quá trình này, các giá trị trung gian $\mathbf{z}^{(l)}$ và $\mathbf{a}^{(l)}$ được lưu lại vì bước lan truyền ngược sẽ cần đến chúng.

Sau đó, \textbf{hàm mất mát} (loss) đo độ lệch giữa dự đoán $\hat{\mathbf{y}}$ và nhãn thật $\mathbf{y}$. Hai lựa chọn phổ biến là
\begin{align}
    \mathcal{L}_{\mathrm{MSE}} &= \frac{1}{N}\sum_{i=1}^{N}\bigl\lVert \hat{\mathbf{y}}_i-\mathbf{y}_i\bigr\rVert_2^{2}
    && \text{(hồi quy)}, \\
    \mathcal{L}_{\mathrm{CE}} &= -\frac{1}{N}\sum_{i=1}^{N}\sum_{k=1}^{K} y_{ik}\log\hat{y}_{ik}
    && \text{(phân loại)}.
\end{align}
Loss là một đại lượng vô hướng: nó cho biết mô hình đang sai bao nhiêu và là mục tiêu cần cực tiểu hóa. Đồng thời, nó là điểm xuất phát của lan truyền ngược.

\subsection{Backpropagation và cập nhật tham số}
Để giảm loss, ta cần biết mỗi tham số ảnh hưởng đến loss như thế nào, tức là \textbf{gradient} $\partial\mathcal{L}/\partial\mathbf{W}^{(l)}$ và $\partial\mathcal{L}/\partial\mathbf{b}^{(l)}$. \textbf{Backpropagation} tính đồng thời gradient của mọi lớp một cách hiệu quả nhờ \textbf{quy tắc chuỗi} (chain rule), đi ngược từ lớp đầu ra về lớp đầu vào và tái sử dụng kết quả trung gian.

Đặt sai số cục bộ $\boldsymbol{\delta}^{(l)}=\partial\mathcal{L}/\partial\mathbf{z}^{(l)}$. Khi đó
\begin{align}
    \boldsymbol{\delta}^{(L)} &= \frac{\partial\mathcal{L}}{\partial\mathbf{a}^{(L)}}\odot\sigma'\!\bigl(\mathbf{z}^{(L)}\bigr), \label{eq:delta-out}\\
    \boldsymbol{\delta}^{(l)} &= \Bigl(\mathbf{W}^{(l+1)}\Bigr)^{\!\top}\boldsymbol{\delta}^{(l+1)}\odot\sigma'\!\bigl(\mathbf{z}^{(l)}\bigr), \label{eq:delta-rec}\\
    \frac{\partial\mathcal{L}}{\partial\mathbf{W}^{(l)}} &= \boldsymbol{\delta}^{(l)}\Bigl(\mathbf{a}^{(l-1)}\Bigr)^{\!\top},
    \qquad
    \frac{\partial\mathcal{L}}{\partial\mathbf{b}^{(l)}} = \boldsymbol{\delta}^{(l)}, \label{eq:grad}
\end{align}
với $\odot$ là tích Hadamard (nhân từng phần tử). Có gradient rồi, tham số được cập nhật theo \textbf{gradient descent} với tốc độ học $\eta>0$:
\begin{equation}
    \mathbf{W}^{(l)}\leftarrow\mathbf{W}^{(l)}-\eta\,\frac{\partial\mathcal{L}}{\partial\mathbf{W}^{(l)}},
    \qquad
    \mathbf{b}^{(l)}\leftarrow\mathbf{b}^{(l)}-\eta\,\frac{\partial\mathcal{L}}{\partial\mathbf{b}^{(l)}}.
    \label{eq:gd}
\end{equation}
Toàn bộ một vòng huấn luyện gồm bốn bước lặp lại: \emph{forward} $\rightarrow$ \emph{tính loss} $\rightarrow$ \emph{backward} $\rightarrow$ \emph{cập nhật tham số}.

\begin{notebox}
Công thức đệ quy~\eqref{eq:delta-rec} cho thấy $\boldsymbol{\delta}^{(l)}$ được tạo ra bằng cách nhân liên tiếp các ma trận $(\mathbf{W}^{(k)})^{\top}$ và các đạo hàm $\sigma'$ qua từng lớp. Một chuỗi tích dài như vậy có thể co lại hoặc phình ra rất nhanh. Đây chính là gốc rễ của hiện tượng vanishing gradient được phân tích ở phần tiếp theo.
\end{notebox}


\section{Dấu hiệu nhận biết Vanishing Gradient}
\label{sec:vanishing}
\textit{Người phụ trách: Trang, Trí}

\subsection{Cơ chế hình thành}
Khai triển đệ quy~\eqref{eq:delta-rec} cho thấy sai số tại lớp $l$ là tích của nhiều thừa số:
\begin{equation}
    \boldsymbol{\delta}^{(l)}
    =\Biggl[\prod_{k=l+1}^{L}\operatorname{diag}\!\bigl(\sigma'(\mathbf{z}^{(k-1)})\bigr)\bigl(\mathbf{W}^{(k)}\bigr)^{\!\top}\Biggr]\boldsymbol{\delta}^{(L)}.
    \label{eq:product}
\end{equation}
Nếu chuẩn của mỗi thừa số nhỏ hơn $1$, gradient co lại gần như theo hàm mũ khi đi ngược về các lớp đầu. Ví dụ, với sigmoid ta có $\sigma'(z)\le 0{,}25$; chỉ riêng thành phần này, sau $10$ lớp gradient đã bị nhân với không quá $0{,}25^{10}\approx 10^{-6}$. Khi đó các lớp gần đầu vào nhận được tín hiệu học gần như bằng không, dù lớp cuối vẫn học bình thường. Hiện tượng này được gọi là \textbf{vanishing gradient}.

\subsection{Các dấu hiệu quan sát được}
Các dấu hiệu dưới đây có thể được kiểm tra trực tiếp trong quá trình huấn luyện; Bảng~\ref{tab:signs} tóm tắt cách phát hiện từng dấu hiệu.

\begin{enumerate}
    \item \textbf{Gradient ở các lớp gần đầu vào rất nhỏ hoặc gần bằng $0$.} Chuẩn gradient $\lVert\nabla_{\mathbf{W}^{(l)}}\mathcal{L}\rVert$ của lớp đầu nhỏ hơn lớp cuối nhiều bậc độ lớn (chẳng hạn $10^{-8}$ so với $10^{-2}$), và giảm đều theo chiều từ lớp cuối về lớp đầu.
    \item \textbf{Trọng số của các lớp đầu thay đổi rất ít qua nhiều epoch.} Bước cập nhật $\eta\,\nabla\mathcal{L}$ gần như bằng $0$ nên các trọng số này gần như đứng yên so với giá trị khởi tạo.
    \item \textbf{Loss giảm rất chậm, sớm tạo plateau hoặc mô hình không hội tụ như kỳ vọng.} Learning curve gần như phẳng dù vẫn còn nhiều epoch và không có dấu hiệu overfitting (loss huấn luyện cũng không giảm).
    \item \textbf{Các activation rơi vào vùng bão hòa của sigmoid/tanh.} Khi $|z|$ lớn, đầu ra tiến sát $0$ hoặc $1$ (sigmoid), $\pm1$ (tanh), khiến $\sigma'(z)\approx 0$ và chặn dòng gradient đi qua. Histogram activation khi đó dồn về hai đầu mút.
    \item \textbf{Lớp cuối học được trong khi các lớp đầu gần như không học.} Mô hình chỉ tinh chỉnh phần đầu ra trên những đặc trưng gần như ngẫu nhiên, nên độ chính xác bị chặn ở mức thấp.
    \item \textbf{Với RNN, mô hình ghi nhớ kém các phụ thuộc dài hạn.} Gradient lan truyền qua nhiều bước thời gian bị co lại, nên mô hình chỉ nắm được ngữ cảnh gần và bỏ sót thông tin ở xa.
\end{enumerate}

\begin{table}[htbp]
    \centering
    \caption{Tóm tắt dấu hiệu vanishing gradient và cách kiểm tra.}
    \label{tab:signs}
    \renewcommand{\arraystretch}{1.25}
    \begin{tabularx}{\textwidth}{@{}>{\raggedright\arraybackslash}p{0.30\textwidth} >{\raggedright\arraybackslash}X >{\raggedright\arraybackslash}X@{}}
        \toprule
        Dấu hiệu & Đại lượng cần theo dõi & Biểu hiện điển hình \\
        \midrule
        Gradient lớp đầu $\approx 0$ & Gradient norm theo từng lớp & Giảm theo cấp số nhân về phía lớp đầu \\
        Trọng số lớp đầu ít đổi & $\lVert\mathbf{W}^{(l)}_t-\mathbf{W}^{(l)}_0\rVert$ theo epoch & Gần như phẳng \\
        Loss giảm chậm, plateau & Learning curve (train loss) & Đường cong phẳng sớm \\
        Activation bão hòa & Histogram của $\mathbf{a}^{(l)}$ & Dồn về $0/1$ hoặc $\pm1$ \\
        Chỉ lớp cuối học & So sánh tốc độ thay đổi giữa các lớp & Lớp cuối thay đổi, lớp đầu đứng yên \\
        RNN quên ngữ cảnh xa & Gradient theo bước thời gian & Suy giảm nhanh khi lùi xa \\
        \bottomrule
    \end{tabularx}
\end{table}

\subsection{Phân biệt với các nguyên nhân khác}
Loss giảm chậm chưa chắc do vanishing gradient. Cần loại trừ các nguyên nhân thường gặp khác như tốc độ học $\eta$ quá nhỏ, lỗi cài đặt (quên gọi \texttt{zero\_grad}, sai nhãn), hay ``dead ReLU'' (neuron luôn xuất ra $0$). Kết luận chỉ đáng tin khi \emph{đồng thời} thấy gradient các lớp đầu nhỏ hơn hẳn các lớp cuối \emph{và} activation bão hòa (hoặc đạo hàm nhỏ) ở những lớp đó.

\section{Cơ chế Vanishing Gradient trong Backpropagation}
Xét một mạng Perceptron đa lớp (MLP) gồm $L$ lớp (độ sâu $L$). Quá trình lan truyền tiến (foward pass) biến đổi vector đầu vào $x_0 \in \mathbb{R}^{N_0}$ thành vector đầu ra $x_L \in \mathbb{R}^{N_L}$ thông qua chuỗi của các lớp ẩn:
\begin{align}
    y_l &= W_l x_l + b_l, \quad l = 0, 1, \dots, L-1 \\
    x_{l+1} &= \sigma(y_l), \quad l = 0, 1, \dots, L-1
\end{align}
trong đó $x_l \in \mathbb{R}^{N_l}$ là vector kích hoạt tại lớp $l$, $W_l \in \mathbb{R}^{N_{l+1} \times N_l}$ là ma trận trọng số, $b_l \in \mathbb{R}^{N_{l+1}}$ là vector độ lệch (bias), và $sigma(\cdot)$ là hàm kích hoạt phi tuyến theo từng phần tử.

Giả sử hàm mất mát trên một mẫy dữ liệu là $\mathcal{l}(y, x_l)$. Đạo hàm của hàm mất mát với ma trận trọng số $W_l$ tại lớp thứ $s$ ($0 \le s \le L-1$) được tính theo thuật toán lan truyền ngược thông qua tích ngoài (outer product):
\begin{equation}
    \frac{\partial \mathcal{l}}{\partial W_s} = \left[ \text{diag}(d_s) B_s \right] \otimes x_s^T
\end{equation}
Trong đó:
\begin{itemize}
    \item $d_x = \sigma'(y_s) \in \mathbb{R}^{N_{s+1}}$ là vector đạo hàm của hàm kích hoạt tại lớp $s$
    \item $\text {diag}(d_s)$ là ma trận đường chéo kích thước $N_{s+1} \times N_{s+1}$ chứa các phần tử của $d_s$.
    \item $E = \frac{\partial \mathcal{L}}{\partial x_L}$ là vector lỗi tại đầu ra (ví dự $E = -2(y - x_L)$ đối với hàm MSE).
    \item $B_s(E)$ là toán tử tuyến tính lan truyền ngược vector lỗi từ lớp $L$ về lớp $s+1$:
    \begin{equation}
        B_s(E)= \left( \frac{\partial x_l}{\partial x_{s+1}} \right)^T E
    \end{equation}
\end{itemize}

% a. Bản chất Toán học của Hiện tượng Vanishing Gradient
\textbf{a. Bản chất Toán học của Hiện tượng Vanishing Gradient}\\

Ma trận Jacobian chuyển tiếp giữa hai lớp liên tiếp $l$ và $l+1$ có dạng:
\begin{equation}
    \frac{\partial x_{l+1}}{\partial x_l} = \text{diag}(d_l) W_l
\end{equation}
Do đó, ma trận Jacobian Đầu vào - Đầu ra (Input - Output Jacobian - IOJ) $\frac{\partial x_L}{\partial x_{S+1}}$ biểu diễn sự lan truyền đạo hàm qua $L-s-1$ lớp:
\begin{equation}
    \frac{\partial x_l}{\partial x_{s+1}} = \frac{\partial x_L}{\partial x_{L-1}} \frac{\partial x_{L-1}}{\partial x_{L-2}} \dots \frac{\partial x_{s+2}}{\partial x_{s+1}} = \prod_{l=s+1}^{L-1} \text{diag}(d_l) W_l
\end{equation}

\textbf{Sự Suy giảm lũy thừa của Tích Ma trận:} Ma trận IOJ toàn cục của mạng $\frac{\partial x_L}{\partial x_1}$ là tích của $L-1$ ma trận:
\begin{equation}
    \frac{\partial x_L}{\partial x_1} = \prod_{l=1}^{L-1} \text{diag}(d_l) W_l
\end{equation}
Giả sử hàm kích hoạt $\sigma(\cdot)$ có đạo hàm bị chặn trên bởi $r > 0$ (ví dụ: $r = 0.25$ đối với Sigmoid, $r = 1$ đối với Tanh/ReLU). Đặt $\Omega_{\max} = \max_{l} \|W_l\|_2$. Theo tính chất sub-multiplicative của chuẩn ma trận:
\begin{equation}
\left\| \frac{\partial x_L}{\partial x_1} \right\|_2 \le \prod_{l=1}^{L-1} \|\text{diag}(d_l)\|_2 \|W_l\|_2 \le \left( r \Omega_{\max} \right)^{L-1}
\end{equation}

\paragraph{Hậu quả:} Khi $r \Omega_{\max} <1$, chuẩn của ma trận IOJ suy giảm theo hàm mũ $\mathcal{O}\left((r \Omega_{\max})^L\right)$ và chiều sâu $L \to \infty$:
\begin{equation}
    \lim_{L \to \infty} \frac{\partial \mathcal{L}}{\partial W_s} = 0 \quad(\text{với các lớp} s \ll L)
\end{equation}
Hiện tượng này khiến cho tín hiệu lỗi $E$ bị triệt tiêu hoàn toàn trước khi tới được các lớp đầu tiên, khiến các trọng số ở xa đầu ra bị "đóng băng" (gradient immobility).

% b. Phân tích theo Lý thuyết Hệ Động lực \& Định lý Furstenberg-Kifer
\textbf{b. Phân tích theo Lý thuyết Hệ Động lực \& Định lý Furstenberg-Kifer}

\textbf{Chuyển đổi DNN sang Hệ Động lực Thời gian Rời rạc} \\
Coi một mạng nơ-ron sâu gồm chuỗi $N$ khối là một hệ động lực phi tuyến rời rạc. Cho các tham số $\omega_n \in \mathbb{R}^s$ là i.i.d. với phân bố $\nu \in \mathcal{P}(\mathbb{R}^s)$. Chuỗi trạng thái được định nghĩa:
\begin{equation}
x_{n+1} = f_{\omega_n}(x_n), \quad n \ge 0
\end{equation}
Sự lan truyền tiến của vector đầu vào $x_0$ tương ứng với quỹ đạo $x_N = f_{\omega_N} \circ f_{\omega_{N-1}} \circ \dots \circ f_{\omega_0}(x_0)$.
Sự tiến hóa của nhiễu loạn tuyến tính $\varepsilon_0 \in \mathbb{R}^d$ qua các lớp tuân theo động lực không dừng:
\begin{equation}
\varepsilon_N = D_N(\omega_N)(x_N) \varepsilon_0 = \left( \prod_{m=0}^N D(\omega_{N-m})(x_{N-m}) \right) \varepsilon_0
\end{equation}
với $D(\omega)(x)$ là ma trận Jacobian $d \times d$ của $f_\omega$ tại điểm $x$. Khi $x_n \to x^*$ (điểm cố định định mệnh), các ma trận $D(\omega_n)(x^*)$ trở thành chuỗi ma trận ngẫu nhiên độc lập cùng phân bố (i.i.d.) có quy luật $\mu \in \mathcal{P}(M)$.

\textbf{Định lý Furstenberg-Kifer và Phổ Số mũ Lyapunov} \\
Xét không gian xạ ảnh $\mathbb{P}^{d-1}$ (tập các đường thẳng đi qua gốc tọa độ trong $\mathbb{R}^d$). Định nghĩa hạt nhân chuyển tiếp Markov $\phi: \mathbb{P}^{d-1} \to \mathcal{P}(\mathbb{P}^{d-1})$:
\begin{equation}
\int_{\mathbb{P}^{d-1}} f(y) \phi(dy|x) = \int_{M} f\left( \frac{A x}{\|A x\|} \right) \mu(dA)
\end{equation}
Tập các độ đo dừng $\mathcal{S} = \{ \varrho \in \mathcal{P}(\mathbb{P}^{d-1}) : \mu \ast \varrho = \varrho \}$ là một hình simplex lồi compact không rỗng.

\begin{theorem}[Furstenberg-Kifer]
Tồn tại dãy các không gian con $\mathbb{R}^d = L_0 \supset L_1 \supset L_2 \dots \supset L_r \supset \{0\}$ và chuỗi các số thực $\lambda_0(\mu) > \lambda_1(\mu) > \dots > \lambda_r(\mu)$ sao cho với $v \in L_i \setminus L_{i+1}$:
\begin{equation}
    \lim_{N \to \infty} \frac{1}{N} \log \|X_N X_{N-1} \dots X_1 v\| = \lambda_i(\mu) \quad \text{a.s.}
\end{equation}
Giá trị của các số mũ Lyapunov $\lambda_i(\mu)$ được xác định bởi:
\begin{equation}
    \lambda(\mu, \varrho) = \int_{M} \int_{\mathbb{P}^{d-1}} \log \|A u\| \mu(dA) \varrho(du)
\end{equation}
khi $\varrho$ biến thiên trên tập độ đo dừng $\mathcal{S}$.
\end{theorem}

\textbf{Tác động ổn định của Kết nối Dư (ResNet):}
Khối kết nối dư biến đổi ánh xạ $x \mapsto f_\omega(x)$ thành $x \mapsto x + f_\omega(x)$, tương ứng với việc thay thế ma trận $A$ bằng $I + A$.

\begin{theorem}[Borkar, 2025]
Phổ số mũ Lyapunov dưới tác động của kết nối dư là một nhiễu loạn nhỏ hơn về phía phổ của ma trận đơn vị $I$ (vector toàn số 1) so với khi không có kết nối dư.
\end{theorem}
\textbf{Tóm tắt chứng minh:}
Xét các vector $\hat{u}_0 = \hat{u}$, $\hat{u}_1 = A \hat{u}$, và $\hat{u}_2 = (I + A)\hat{u} = \hat{u}_0 + \hat{u}_1$ với $\|\hat{u}\| = 1$. Vector $\hat{u}_3$ là giao điểm các đường chéo của hình bình hành tạo bởi $0, \hat{u}_0, \hat{u}_1, \hat{u}_2$. Khi chiếu lên mặt cầu $S^d$ và không gian xạ ảnh $\mathbb{P}^{d-1}$, khoảng cách metric giữa $\pi(\hat{u}_3)$ và $\pi(\hat{u}_0)$ nhỏ hơn khoảng cách giữa $\pi(\hat{u}_1)$ và $\pi(\hat{u}_0)$.Điều này dẫn đến:
\begin{equation}
\int_{\mathbb{P}^{d-1}} \int_{M} \log \|(I+A)u\| \mu(dA) \varrho(du) - \int_{\mathbb{P}^{d-1}} \int_{M} \log \|Au\| \mu(dA) \varrho(du) < 0
\end{equation}
Kết quả này chứng minh rằng kết nối dư kéo các số mũ Lyapunov tiệm cận về 0 (tương ứng chuẩn bằng 1), triệt tiêu sự bùng nổ hoặc suy giảm theo hàm mũ.


% c. Giải pháp Bộ lọc Cộng Đẳng hướng Ngẫu nhiên (roaFNN \& roaRNN)
\textbf{c. Giải pháp Bộ lọc Cộng Đẳng hướng Ngẫu nhiên (roaFNN \& roaRNN)}

\textbf{Kiến trúc roaFNN và Động lực học Tích số}\\
Mô hình roaFNN thiết lập một tổ hợp lồi giữa biến đổi phi tuyến tiêu chuẩn và ma trận bán đẳng hướng ngẫu nhiên $O_l$:
\begin{equation}
x_{l+1} = \alpha \sigma(W_l x_l + b_l) + (1-\alpha) O_l x_l, \quad l = 0, \dots, L-1
\end{equation}
trong đó $O_l \in \mathbb{R}^{N_{l+1} \times N_l}$ thỏa mãn $O_l^T O_l = I_{N_l}$ (nếu $N_l \le N_{l+1}$) hoặc $O_l O_l^T = I_{N_{l+1}}$ (nếu $N_l \ge N_{l+1}$). Ma trận Jacobian từng lớp trở thành:
\begin{equation}
\frac{\partial x_{l+1}}{\partial x_l} = \alpha \text{diag}(d_l) W_l + (1-\alpha) O_l
\end{equation}

\textbf{ Các Bổ đề Biên và Định lý Độc lập với Chiều sâu}

\begin{lemma}[Biên trên từng lớp]
Với $P_l = \frac{\partial x_L}{\partial x_l}$, ta có: $\|P_l\| \le \left[ 1 + \alpha (r \Omega_l - 1) \right] \|P_{l+1}\|$.
\end{lemma}

\begin{lemma}[Biên dưới từng lớp]
Giả sử dãy kích thước lớp không tăng ($N_0 \ge N_1 \ge \dots \ge N_L$). Nếu $\alpha < \frac{1}{1 + r \Omega_{\max}}$, ta có: $\|P_l\| \ge \left[ 1 - \alpha (1 + r \Omega_l) \right] \|P_{l+1}\|$.
\end{lemma}

\begin{theorem}[Ceni, 2022 - Biên Jacobian độc lập với độ sâu $L$]
Xét mô hình roaFNN với tham số hóa $\alpha = \rho (L-1)^{-1}$ ($\rho > 0$). Giá trị trị riêng lớn nhất (chuẩn) của ma trận IOJ $\frac{\partial x_L}{\partial x_1}$ luôn thỏa mãn biên trên độc lập với $L$:
\begin{equation}
\left\| \frac{\partial x_L}{\partial x_1} \right\|_2 \le \exp\left( \rho (r \Omega_{\max} - 1) \right)
\end{equation}
Đồng thời, nếu $\rho < \frac{L-1}{1 + r \Omega_{\max}}$ và kích thước các lớp không tăng, ta có biên dưới độc lập với $L$:
\begin{equation}
\exp\left( -\rho (1 + r \Omega_{\max}) \right) \le \left\| \frac{\partial x_L}{\partial x_1} \right\|_2
\end{equation}
\end{theorem}

\textbf{Nhận xét:} Khi $L \to \infty$, cả hai biên trên và biên dưới đều hội tụ về các hằng số dương hữu hạn hoàn toàn không phụ thuộc vào độ sâu $L$. Điều này chứng minh toán học rằng roaFNN triệt tiêu hoàn toàn hiện tượng Vanishing Gradient lẫn Exploding Gradient mà không cần Batch Normalization hay Gradient Clipping.

% d. Mở rộng cho Mạng Hồi quy (roaRNN) \& Mô hình Thời gian Liên tục
\textbf{d. Mở rộng cho Mạng Hồi quy (roaRNN) \& Mô hình Thời gian Liên tục}
Mô hình roaRNN áp dụng ma trận đẳng hướng $O$ dùng chung qua các bước thời gian $k$:
\begin{equation}
x_{k+1} = \alpha \sigma(W_h x_k + b_h + W_i u_{k+1}) + (1-\alpha) O x_k
\end{equation}
Trong góc nhìn thời gian liên tục, roaRNN chính là dạng xấp xỉ Euler hiện của phương trình vi phân thông thường (ODE):
\begin{equation}
\tau \dot{x} = A x + \sigma(W_h x + b_h + W_i u)
\end{equation}
với hằng số thời gian $\tau = \rho^{-1}$ và ma trận ma sát/dẫn dắt $A = \alpha^{-1} \left[ (1-\alpha)O - I \right]$.

\textbf{Bảng So sánh Tổng hợp các Phương pháp Ổn định Đạo hàm}

\begin{table}[h]
\centering
\small
\begin{tabular}{p{3.5cm}p{3.5cm}p{4cm}p{4cm}}
\toprule
\textbf{Phương pháp} & \textbf{Biểu diễn Toán học} & \textbf{Cơ chế Động lực học} & \textbf{Giới hạn / Nhược điểm} \\
\midrule
\textbf{Vanilla MLP / RNN} & $x_{l+1} = \sigma(W_l x_l + b_l)$ & Tích Jacobian $\prod \text{diag}(d_l)W_l$ suy giảm/bùng nổ lũy thừa & Bị triệt tiêu đạo hàm nghiêm trọng khi $L > 20$ hoặc $T > 100$. \\
\addlinespace
\textbf{ResNet} & $x_{l+1} = x_l + f_\omega(x_l)$ & Thu hẹp phổ Lyapunov quanh $I$ ($A \mapsto I + A$) & Vẫn phụ thuộc vào Batch Normalization ở độ sâu cực lớn. \\
\addlinespace
\textbf{roaFNN / roaRNN} & $x_{l+1} = \alpha \sigma(\cdot) + (1-\alpha)O_l x_l$ & Chặn phổ giá trị đơn độc lập độ sâu $L$ qua $\alpha = \frac{\rho}{L-1}$ & Chi phí bộ nhớ tăng nhẹ do lưu trữ ma trận $O_l$. \\
\bottomrule
\end{tabular}
\caption{So sánh đặc tính toán học của các kiến trúc mạng nơ-ron sâu.}
\end{table}


\subsection{Chain rule trong mạng nhiều lớp}
Với mạng sâu, gradient tại một lớp sớm là tích của nhiều đạo hàm cục bộ:
\[
\frac{\partial \mathcal{L}}{\partial \mathbf{h}^{(l)}}
=
\left(\prod_{k=l+1}^{L}
\frac{\partial \mathbf{h}^{(k)}}{\partial \mathbf{h}^{(k-1)}}\right)
\frac{\partial \mathcal{L}}{\partial \mathbf{h}^{(L)}}.
\]

\subsection{Sự co nhỏ qua phép nhân lặp}
Nếu chuẩn của các Jacobian thường nhỏ hơn $1$, tích lặp qua nhiều lớp sẽ giảm theo cấp số nhân, khiến gradient truyền về các lớp đầu trở nên rất nhỏ.

\subsection{Liên hệ giữa đạo hàm activation và trọng số}
Phân tích cách $f'(z)$ và phổ của ma trận trọng số cùng quyết định độ lớn gradient.

\subsection{Ví dụ số tối giản}
Đề xuất tính tay một chuỗi gồm nhiều đạo hàm bằng $0.25$ để cho thấy $0.25^L$ nhanh chóng tiến về $0$ khi $L$ tăng.

\section{Các yếu tố gây Vanishing Gradient}

\subsection{Hàm activation}
% \textit{Người phụ trách: Khanh}
% \begin{itemize}
%     \item Sigmoid và tanh có vùng bão hòa với đạo hàm rất nhỏ.
%     \item Sigmoid có đạo hàm cực đại chỉ bằng $0.25$.
%     \item So sánh với ReLU và các biến thể để làm rõ ảnh hưởng của activation.
% \end{itemize}
Trong ma trận Jacobian từng lớp $\frac{\partial x_{l+1}}{\partial x_l} = \text{diag}(d_l) W_l$, ma trận đường chéo $\text{diag}(d_l)$ chứa đạo hàm của hàm kích hoat $d_l = \sigma'(y_l)$ đóng vai trò là hệ số điều tiết biên độ gradient khi truyền qua mội lớp.

Giả sử đạo hàm của hàm kích hoạt bị chặn trên bởi hằng số $r > 0$:
\begin{equation}
    r = \sup_{y \in \mathbb{R}} |\sigma'(y)| \implies \|\text{diag} (d_l)\|_2 \le r
\end{equation}

Nếu $r < 1$, ma trận $\text{diag}(d_l)$ đóng vai trò là một toán tử thu nhỏ )contraction operator), trực tiếp làm suy giảm biên độ tín hiệu lỗi qua từng lop81 lan truyền ngược.

% a.Hàm Sigmoid (Logistic Function)
\textbf{a. Hàm Sigmoid (Logistic Function)}
\begin{equation}
\sigma(y) = \frac{1}{1 + e^{-y}}
\end{equation}
Đạo hàm cấp một của hàm Sigmoid có dạng:
\begin{equation}
 \sigma'(y) = \sigma(y)\left(1 - \sigma(y)\right)
\end{equation}

\begin{itemize}
    \item \textbf{Giá trị cực đại:} $sigma'(y)$ đạt giá trị lớn nhất tại $y = 0$ với $r = \max_y \sigma'(y) = 0.25$.
    \item \textbf{Tác động đến Gradient:}: Vì $\|\text{diag}(d_l)\|_2 \le 0.25$, ngay cả khi ma trận trọng số thỏa mãn $ \|W_l\|_2 = 1$, tín hiệu gradient khi đi qua $L-1$ lớp Sigmoid sẽ bị thu nhỏ với tốc độ ít nhất là":
    \begin{equation}
        \left \| \prod_{l=1}^{L-1} \text{diag}(d_l) \right\|_2 \le (0.25)^{L-1} = 4^{-(L-1)}
    \end{equation}
    \item \textbf{Hiện tượng Bão hòa (Saturation) :} Khi $|y| \gg 0$ (nguồn tín hiệu đầu vào lớp quá lớn hoặc quá nhỏ), $\sigma(y) \to 1$ hoặc $\sigma(y) \ to 0$, dẫn đến $\sigma(y) \to 0$, dẫn đến $\sigma'(y) \to 0$. Vùng bão hòa này khiến cho $d_l \approx 0$, lập tức triệt tiêu gradient truyền qua lớp đó.
\end{itemize}

% b. Hàm tangent Hyperbol (Tanh):
\textbf{b. Hàm tangent Hyperbol (Tanh)}\\
Hàm Tanh ép giá trị đầu ra về khoảng $(-1, 1)$:
\begin{equation}
    \sigma(y) = \tanh(y) = \frac{e^y - e^{-y}}{e^y + e^{-y}}
\end{equation}
\text{Đạo hàm của hàm Tanh được tính theo công thức:}
\begin{equation}
    \sigma'(y) = 1 - \tanh^2(y)
\end{equation}

\begin{itemize}
    \item \textbf{Giá trị cực đại:} $r = \max_y \sigma'(y) = 1.0$ (đạt được tại $y = 0$).
    \item \textbf{So sánh với Sigmoid:} Tanh cải thiện đáng kể so với Sigmoid vì $r = 1.0$ giúp gradient không bị thu nhỏ cưỡng bức tại vùng làm việc tuyến tính ($y \approx 0 $)
    \item \textbf{Hạn chế bão hòa:} Tuy nhiên, khi $|y| > 2$, $\tanh^2(y) \to 1$ khiến $\sigma'(y) \to 0$. Do đó , nếu trọng số $w_l$ khiến giá trị $y_l$ rơi vào vùng bão hòa hai đầu, hiện tượng Vanishing Gradient vẫn xảy ra nghiêm trọng.
\end{itemize}

% c. Hàm Tuyến tính Từng phần (ReLU \& Leaky ReLU):
\textbf{c. Hàm Tuyến tính Từng phần (ReLU \& Leaky ReLU)}\\
Để khắc phục hiện tượng bão hòa đạo hàm của các hàm dạng Sigmoid và Tanh, hàm ReLu (Retified Linear Unit) được sử dụng:
\begin{equation}
    \sigma(y) = \max(o, y) \implies \sigma'(y) = \begin{cases} 1 & \text{nếu } y > 0 \\ 0 & \text{nếu } y \le 0 \end{cases}
\end{equation}
\begin{itemize}
    \item \textbf{Ưu điểm giải quyết Vanishing Gradient:} Đối với các neuron ở trạng thai 1kích hoạt ($y > 0$), đạo hàm luôn bằng $1$ ($r = 1.0$). Điều này loại bỏ hoàn toàn hệ số thu nhỏ đạo hàm của hàm kích hoạt, giúp gradient truyền qua nhiều lop81 mà không bị suy giảm do $\sigma'(y)$.
    \item \textbf{Vấn đề "Dying ReLU":} Khi $y \le 0$, $sigma'(y) = 0$. Nếu một neuron bị rơi vào vùng âm với trọng số cập nhật không phù hợp, nó sẽ vĩnh viễn không phát tín hiệu gradient về sau ($\text{diag}(d_l) = 0$).
    \item \textbf{Biến thể Leaky ReLu/ Parametric ReLU:} Khắc phục điểm chế bằng cách gán một độ dốc nhỏ $\alpha \approx 0.01$ cho vùng âm:
    \begin{equation}
        \sigma'(y) = \begin{cases} 1 & \text{nếu } y > 0 \\  \alpha & \text{nếu } y \le 0 \end{cases}
    \end{equation}
    Giúp duy trì dòng chảy gradient $\alpha >0 $ liên tục qua mọi trạng thái đầu vào.
\end{itemize}

% d.Sự Tương tác Giữa Đạo hàm Activation và Chuẩn Trọng số
\textbf{d. Sự Tương tác Giữa Đạo hàm Activation và Chuẩn Trọng số:} Độ lớn tổng thể của ma trận Jacobian từng lớp $J_l = \frac{\partial x+{l+1}}{\partial x_l}$ chịu ảnh hưởng bởi tích giữa chuẩn đạo hàm activation và chuẩn ma trận trọng số:
\begin{equation}
    \|J_l\|_2 = \|\text{diag}(d_l) W_l\|_2 \le \|\text{diag}(d_l)\|_2 \|W_l\|_2 \le r \|W_l\|_2
\end{equation}

\begin{itemize}
    \item Nếu $r \|W_l\|_2 < 1$, ma trận $J_l$ là toán tử co, dẫn đến vanishing gradient khi chuỗi tích $\prod_{l=1}^{L-1} J_l$ tiến về $0$.
    \item Nếu $r \|W_l\|_2 > 1$, tích ma trận có nguy cơ bùng nổ gradient (exploding gradient).
    \item Để đạt trạng thái cân bằng \textit{Dynamical Isometry} ($\|J_l\|_2 \approx 1$), việc lựa chọn hàm kích hoạt có $r = 1$ (như Tanh hoặc ReLU) phải đi kèm với các kỹ thuật khởi tạo trọng số phù hợp (như Xavier/Kaiming Initialization) để đảm bảo $\|W_l\|_2 \approx 1$.
\end{itemize}

\subsection{Khởi tạo và độ lớn của weight}
\label{subsec:weight_init}

Huấn luyện mạng neural là một quá trình lặp: tham số được cập nhật dần từng bước, bắt đầu từ một giá trị xuất phát nào đó. Do đó, trước bước cập nhật đầu tiên, ta buộc phải chọn giá trị ban đầu cho toàn bộ trọng số. Hàm mất mát của một mạng sâu có hình dạng phức tạp và không lồi, nên kết quả huấn luyện phụ thuộc đáng kể vào lựa chọn này~\cite[\S8.4]{long-goodfellow2016}.

Ảnh hưởng của điểm xuất phát thể hiện ở ba khía cạnh. Thứ nhất, nó quyết định quá trình huấn luyện có chạy được hay không: một số giá trị ban đầu gây ra lỗi số học và làm thuật toán thất bại ngay từ đầu. Thứ hai, nó ảnh hưởng đến tốc độ hội tụ. Thứ ba, ngay cả khi quá trình học hội tụ, nó vẫn có thể quyết định mô hình dừng ở một điểm có loss thấp hay cao~\cite[\S8.4]{long-goodfellow2016}.

Mục IV đã chỉ ra rằng gradient tại một lớp ở phía đầu mạng là tích của các ma trận Jacobian cục bộ:
\[ \frac{\partial \mathcal{L}}{\partial \mathbf{h}^{(l)}} = \left( \prod_{k=l+1}^{L} \frac{\partial h^{(k)}}{\partial \mathbf{h}^{(k-1)}} \right) \frac{\partial \mathcal{L}}{\partial \mathbf{h}^{(L)}}, \qquad \frac{\partial \mathbf{h}^{(k)}}{\partial \mathbf{h}^{(k-1)}} = \operatorname{diag}\big( f'(\mathbf{z}^{(k)}) \big) W^{(k)}. \]
Trước khi huấn luyện, các ma trận trọng số $W^{(k)}$ được khởi tạo ngẫu nhiên, do đó độ lớn ban đầu của trọng số quyết định trực tiếp kích thước của từng thừa số trong tích này. Nó cũng quyết định $\mathbf{z}^{(k)}$ vì $\mathbf{z}^{(k)}$ được tính bằng cách truyền đầu vào qua tất cả các lớp ngẫu nhiên trước đó, và vì vậy nó cũng quyết định $f'(\mathbf{z}^{(k)})$ sẽ lớn hay tiến gần về 0.

Schoenholz và cộng sự (2017) ~\cite[\S4]{long-schoenholz2017} cho thấy rằng trên trung bình, mỗi lớp nhân kích thước bình phương trung bình (mean squared size) của tín hiệu lan truyền ngược (backward signal) với một số hệ số khuếch đại phụ thuộc vào phương sai của trọng số và hàm kích hoạt:

\[  \chi \;=\; n \operatorname{Var}(W)\; \mathbb{E}\big[ f'(\mathbf{z})^2 \big], \]
trong đó $n$ là chiều rộng của lớp. Qua $L - l$ lớp, kích thước bình phương trung bình của gradient được nhân lên khoảng $\chi^{L - l}$, tức là độ lớn của nó được nhân lên khoảng $\chi^{(L - l)/2}$ lần. Hiện tượng xảy ra dựa vào giá trị của $\chi$ ~\cite[\S4]{long-schoenholz2017}:
\begin{itemize}
    \item Nếu $\chi < 1$, gradient sẽ bị biến mất (gradient vanishing). Nếu chuỗi mạng đủ sâu (tức $L - l$ lớn), giá trị $\chi^{L-l}$ sẽ tiến rất nhanh về 0. Điều này làm cho gradient tại các lớp đầu tiên ($\frac{\partial \mathcal{L}}{\partial \mathbf{h}^{(l)}}$) cực kỳ nhỏ, khiến các trọng số ở các lớp này gần như không được cập nhật trong quá trình huấn luyện.
    \item Nếu $\chi = 1$, gradient được bảo toàn ổn định qua các lớp
    \item Nếu $\chi > 1$, gradient bùng nổ và tiến tới vô cùng (gradient exploding).
\end{itemize}
Việc khởi tạo sẽ chọn $\operatorname{Var}(W)$, và do đó sẽ quyết định giá trị của $\chi$.

Khi trọng số quá nhỏ, ta sẽ có $\chi < 1$ vì mỗi lớp sẽ làm thu nhỏ cả tín hiệu lan truyền tiến (forward signal) lẫn gradient lan truyền ngược. Ví dụ, với chiều rộng lớp $n = 100$ và trọng số được lấy từ phân phối $\mathcal{N}(0, 0.01^2)$, ta có $n \operatorname{Var}(W) = 10^{-2}$, do đó độ lớn của tín hiệu và gradient bị chia cho 10 ở mỗi lớp, và sau 20 lớp, nó bị nhân với $10^{-20}$. Các lớp ở sâu khi đó sẽ nhận được giá trị kích hoạt rất nhỏ gần như giống hệt nhau với mọi đầu vào (xem Mục V.4), và các lớp ở phía trước sẽ gần như không nhận được gradient nào.

Ngược lại, trọng số quá lớn sẽ làm cho $|\mathbf{z}^{(k)}|$ tăng lên và đẩy các đơn vị vào vùng bão hòa của hàm sigmoid và tanh, nơi $f'(\mathbf{z}^{(k)})$ gần bằng 0 (Mục V.1). Một đơn vị bị bão hòa hầu như không truyền được gradient. Glorot và Bengio (2010) ~\cite{glorot2010} đã quan sát thấy hiện tượng bão hòa này ở các lớp ẩn của mạng sigmoid và tanh sâu.

Tuy nhiên, hiện tượng bão hòa không phải là toàn bộ vấn đề. Các trọng số lớn cũng khuếch đại tín hiệu. Poole (2016)~\cite{poole2016} và Schoenholz (2017) ~\cite{long-schoenholz2017} cho thấy rằng khi phương sai của trọng số vượt qua một giá trị tới hạn (critical value), gradient trung bình sẽ tăng theo độ sâu thay vì thu nhỏ lại. Vì vậy, trọng số quá lớn không chỉ gây bão hòa các unit làm mất gradient, mà còn làm mất ổn định hệ thống và có thể khiến gradient bùng nổ.

Hai quy tắc khởi tạo tiêu chuẩn của Xavier ~\cite{glorot2010} và He ~\cite{he2015rectifiers} chọn $\operatorname{Var}(W)$ sao cho $\chi$ xấp xỉ $1$. Chúng sẽ được thảo luận chi tiết trong Mục \ref{subsec:fix-weightinit-normalization}.

Chúng ta cần lưu ý rằng việc khởi tạo chỉ đảm bảo một điểm bắt đầu tốt. Các phân tích kể trên đều dựa trên giả định rằng các lớp có chiều rộng lớn, các trọng số độc lập với nhau, và nó chỉ mô tả hành vi trung bình tại bước huấn luyện đầu tiên. Đối với các mạng có chiều rộng hữu hạn, Hanin (2018) ~\cite{hanin2018} đã chỉ ra rằng độ biến thiên của gradient giữa các lần khởi tạo ngẫu nhiên sẽ tăng theo tổng các nghịch đảo chiều rộng của các lớp; do đó, một mạng sâu và hẹp vẫn có thể gặp phải hiện tượng triệt tiêu hoặc bùng nổ gradient (Mục \ref{subsec:depth}). Chuẩn hóa theo lô (Batch normalization) và các kết nối tắt (residual connections) giúp quá trình huấn luyện ít bị nhạy cảm hơn đối với quy mô khởi tạo ban đầu. Hai phương pháp này sẽ được trình bày ở Mục \ref{subsec:fix-weightinit-normalization} và VII.3.

\subsection{Độ sâu của mạng}
\label{subsec:depth}

Mục IV.1 đã viết gradient tại lớp $l$ dưới dạng:
\[
\frac{\partial \mathcal{L}}{\partial \mathbf{h}^{(l)}}
= \left( \prod_{k=l+1}^{L} \frac{\partial \mathbf{h}^{(k)}}{\partial \mathbf{h}^{(k-1)}} \right) \frac{\partial \mathcal{L}}{\partial \mathbf{h}^{(L)}}.
\]
Tích này có $L - l$ nhân tử, vì vậy nó trở nên dài hơn khi mạng càng sâu, và nó dài nhất đối với các lớp nằm gằn đầu vào nhất.  Lớp 1 được truyền $L - 1$ ma trận Jacobian, trong khi lớp cuối không phải truyền qua ma trận Jacobian nào. Vì vậy, việc thêm một lớp trong mạng sẽ thêm một nhân tử vào gradient của \textit{mọi} lớp ở phía trước.

Nếu kích thước của mỗi nhân tử đều nhỏ hơn 1, tích này sẽ suy giảm theo cấp số nhân. Nếu $\big\| \partial \mathbf{h}^{(k)} / \partial \mathbf{h}^{(k-1)} \big\| \le \gamma < 1$ với mọi lớp, thì

\[
\left\| \frac{\partial \mathcal{L}}{\partial \mathbf{h}^{(l)}} \right\| \;\le\; \gamma^{\,L-l}\, \left\| \frac{\partial \mathcal{L}}{\partial \mathbf{h}^{(L)}} \right\| .
\]
Độ sâu chính là số mũ trong chặn trên này. Một hệ số $\gamma = 0.9$ trên mỗi lớp là vô hại đối với một mạng nông (ví dụ: $0.9^{10} \approx 0.35$), nhưng nó lại rất nghiêm trọng với một mạng sâu (ví dụ: $0.9^{100}\approx 2.7\times10^{-5}$). Lập luận tương tự theo chiều ngược lại sẽ giải thích hiện tượng bùng nổ gradient (ví dụ: $1.1^{10} \approx 2.6$ và $1.1^{100} \approx 1.4 \times 10^4$). Ví dụ về hàm sigmoid ở Mục IV.4, nơi mỗi nhân tử tối đa là 0.25, là trường hợp cực đoan của chặn trên này.

Mục V.2 đã chỉ ra rằng, trên trung bình, mỗi lớp sẽ scale gradient theo một số hệ số khuếch đại $\chi = n \operatorname{Var}(W) \mathbb{E}[f'(z)^2]$, do đó
\[
\left\| \frac{\partial \mathcal{L}}{\partial \mathbf{h}^{(l)}} \right\| \;\sim\; \chi^{(L-l)/2}.
\]
Hệ số khuếch đại $\chi$ không phụ thuộc vào độ sâu; độ sâu chỉ quyết định hệ số này được áp dụng bao nhiêu lần. Chính vì thế, một hệ số khuếch đại chỉ nhỏ hơn 1 một chút là vô hại trong một mạng nông nhưng lại gây hại trong một mạng sâu. Schoenholz và cộng sự (2017) ~\cite{long-schoenholz2017} đã cho thấy rằng các mạng sâu chỉ có thể huấn luyện được khi việc khởi tạo ở rất gần điểm tới hạn $\chi = 1$, và khởi tạo càng gần điểm này thì mạng càng có thể xây dựng sâu hơn.

\subsection{Sự suy giảm tín hiệu trong Forward Pass}
\label{sec:long-forward}
\textit{Người phụ trách: Long}

Sự lan truyền gradient chịu ảnh hưởng trực tiếp của những trạng thái mà mạng tạo ra trong lượt truyền xuôi (\textit{forward pass}). Nếu các phép biến đổi liên tiếp làm suy yếu độ nhạy đối với đầu vào, tín hiệu sai số truyền ngược cũng có thể bị co nhỏ. Tuy nhiên, ta cần phân biệt ba đại lượng: độ lớn của activation, mức độ khác biệt giữa các biểu diễn và độ lớn của gradient. Activation nhỏ không tự nó là một điều kiện đủ để kết luận rằng đã có vanishing gradient; ngược lại, activation vẫn có thể khác không trong khi độ nhạy và gradient đã suy giảm nghiêm trọng.
\subsubsection{Mô hình truyền tín hiệu và độ nhạy cục bộ}
Xét mạng truyền thẳng gồm $L$ lớp, với
\begin{equation}
 \mathbf h^{(0)}=\mathbf x,\qquad
 \mathbf z^{(\ell)}=\mathbf W^{(\ell)}\mathbf h^{(\ell-1)}+\mathbf b^{(\ell)},\qquad
 \mathbf h^{(\ell)}=\phi\!\left(\mathbf z^{(\ell)}\right).
 \label{eq:long-ff}
\end{equation}
Ở đây, $\mathbf h^{(\ell)}\in\mathbb R^{n_\ell}$ là vector activation, $\mathbf z^{(\ell)}$ là vector tiền kích hoạt, và $\phi$ được áp dụng theo từng phần tử. Trong hai mục này, gradient được biểu diễn bằng vector cột; $\|\cdot\|_2$ là chuẩn Euclid đối với vector và chuẩn phổ đối với ma trận. Jacobian của lớp $\ell$ là
\begin{equation}
 \mathbf J_\ell
 =\frac{\partial\mathbf h^{(\ell)}}{\partial\mathbf h^{(\ell-1)}}
 =\mathbf D_\ell\mathbf W^{(\ell)},\qquad
 \mathbf D_\ell=\operatorname{diag}\!\left(\phi'(\mathbf z^{(\ell)})\right).
 \label{eq:long-jacobian}
\end{equation}
Với nhiễu đầu vào đủ nhỏ $\Delta\mathbf h^{(k)}$, đặt
$\mathbf P_{L:k}=\mathbf J_L\mathbf J_{L-1}\cdots\mathbf J_{k+1}$. Khi các đạo hàm tồn tại, khai triển bậc nhất và quy tắc chuỗi sẽ cho ta
\begin{equation}
 \Delta\mathbf h^{(L)}\approx\mathbf P_{L:k}\Delta\mathbf h^{(k)},\qquad
 \mathbf g^{(k)}=\mathbf P_{L:k}^{\mathsf T}\mathbf g^{(L)},\qquad
 \mathbf g^{(\ell)}=\nabla_{\mathbf h^{(\ell)}}\mathcal L.
 \label{eq:long-duality}
\end{equation}
Như vậy, lan truyền xuôi một nhiễu nhỏ và lan truyền ngược một gradient sẽ sử dụng cùng một tích Jacobian, và sẽ khác nhau bởi phép chuyển vị. Nếu $\|\mathbf P_{L:k}\|_2$ nhỏ thì cả độ nhạy đầu ra đối với nhiễu ở lớp $k$ lẫn gradient truyền về lớp đó đều bị chặn bởi một hệ số nhỏ. Trong trường hợp tích Jacobian có các giá trị kỳ dị rất khác nhau, chỉ một số hướng có thể bị co mạnh; do đó, không thể suy ra hành vi của mọi hướng chỉ từ một thống kê trung bình.

\subsubsection{Suy giảm độ lớn và sự bão hòa của activation}
Một cơ chế suy giảm dễ quan sát là sự co nhỏ của tín hiệu qua nhiều lớp. Trong mô hình tuyến tính đơn giản $\mathbf b^{(\ell)}=\mathbf0$ và $\phi(z)=z$, giả sử các trọng số độc lập, có trung bình bằng không, phương sai $v_\ell$, và độc lập với đầu vào của lớp. Khi các thành phần đầu vào có cùng một moment bậc hai là $q_{\ell-1}$, ta sẽ có
\begin{equation}
 q_\ell=\mathbb E\!\left[(h_i^{(\ell)})^2\right]
 =n_{\ell-1}v_\ell q_{\ell-1}.
 \label{eq:long-linear-variance}
\end{equation}
Kỳ vọng ở đây xét trên khởi tạo ngẫu nhiên và phân phối đầu vào nếu đầu vào cũng được xem là ngẫu nhiên. Khi activation có trung bình bằng không, $q_\ell$ đồng thời là phương sai. Nếu $n_{\ell-1}v_\ell=\alpha<1$ ở mọi lớp thì $q_\ell=\alpha^\ell q_0$: độ lớn bình phương trung bình của tín hiệu sẽ giảm theo cấp số nhân. Đây là mô hình lý tưởng hóa cho lập luận bảo toàn phương sai trong lúc khởi tạo; các giả thiết tuyến tính và độc lập sẽ không còn đúng hoàn toàn trong mạng phi tuyến sau khi huấn luyện~\cite[\S4.2]{glorot2010}.

Một cơ chế khác xuất hiện khi tiền kích hoạt đi vào vùng bão hòa. Với $\phi(z)=\tanh(z)$, nếu $|z|$ lớn thì $|\phi(z)|\approx1$ nhưng $\phi'(z)=1-\tanh^2(z)\approx0$. Khi đó, activation không nhỏ, song một thay đổi nhỏ của đầu vào gần như không làm thay đổi đầu ra neuron. Sự bão hòa làm các phần tử tương ứng của $\mathbf D_\ell$ nhỏ đi, góp phần co Jacobian trong~\eqref{eq:long-jacobian}. Ngược lại, trong lân cận $z=0$, $\tanh(z)\approx z$ và $\tanh'(z)\approx1$, tức activation nhỏ tại đây không đồng nghĩa với việc đạo hàm nhỏ. Các quan sát này cho thấy rằng ta cần phải đánh giá phân bố preactivation và đạo hàm, thay vì chỉ đơn thuần nhìn vào độ lớn của activation~\cite[\S3]{glorot2010}.

Ngoài gradient trạng thái, tín hiệu lan truyền xuôi còn tác động trực tiếp đến gradient trọng số. Đặt $\boldsymbol\delta^{(\ell)}=\nabla_{\mathbf z^{(\ell)}}\mathcal L$. Đối với một mẫu dữ liệu, ta sẽ có
\begin{equation}
 \nabla_{\mathbf W^{(\ell)}}\mathcal L
 =\boldsymbol\delta^{(\ell)}\big(\mathbf h^{(\ell-1)}\big)^{\mathsf T},\qquad
 \left\|\nabla_{\mathbf W^{(\ell)}}\mathcal L\right\|_F
 =\|\boldsymbol\delta^{(\ell)}\|_2\,\|\mathbf h^{(\ell-1)}\|_2.
 \label{eq:long-weight-gradient}
\end{equation}
Do đó, nếu một thừa số suy giảm trong khi thừa số còn lại bị chặn, gradient trọng số cũng sẽ nhỏ đi. Tuy nhiên, hai thừa số còn có thể biến đổi ngược chiều và bù trừ nhau về độ lớn; đối với một batch nhỏ, gradient còn là tổng hoặc là trung bình của nhiều tích ngoài. Chính vì thế, gradient trọng số không thể được dùng để thay thế hoàn toàn cho gradient trạng thái khi ta chẩn đoán về sự lan truyền tín hiệu~\cite[\S4.2--4.3]{glorot2010}.

\subsubsection{Mất khả năng phân biệt đầu vào dù phương sai còn được duy trì}
Sự bảo toàn về độ lớn của từng biểu diễn vẫn là chưa đủ để bảo toàn sự khác biệt giữa các đầu vào với nhau. Để mô tả hiện tượng này, Schoenholz và cộng sự đã xét một họ mạng truyền thẳng rộng, chưa được huấn luyện, có trọng số và bias được khởi tạo một cách độc lập theo
\begin{equation}
 W_{ij}^{(\ell)}\sim\mathcal N\!\left(0,\frac{\sigma_w^2}{n_{\ell-1}}\right),\qquad
 b_i^{(\ell)}\sim\mathcal N(0,\sigma_b^2).
 \label{eq:long-mf-init}
\end{equation}
Trong xấp xỉ trường trung bình (\textit{mean-field approximation}), tiền kích hoạt tại một lớp được thay bằng biến Gaussian có cùng hai moment đầu. Với một đầu vào cố định, đặt $q_\ell=\mathbb E[(z_i^{(\ell)})^2]$; kỳ vọng lúc này lấy trên họ mạng khởi tạo ngẫu nhiên, không phải phương sai của một neuron trên tập dữ liệu. Khi tiền kích hoạt có trung bình bằng không, $q_\ell$ chính là phương sai và thỏa mãn gần đúng
\begin{equation}
 q_\ell\approx\sigma_w^2\int\mathcal Dz\,
 \phi\!\left(\sqrt{q_{\ell-1}}z\right)^2+\sigma_b^2,
 \qquad \mathcal Dz=\frac{e^{-z^2/2}}{\sqrt{2\pi}}\,dz.
 \label{eq:long-mf-q}
\end{equation}
Đây là cách viết lại quan hệ truyền phương sai của~\cite[\S2]{long-schoenholz2017} theo ký hiệu của báo cáo.

Xét hai đầu vào $\mathbf x_a,\mathbf x_b$ đi qua \emph{cùng một mạng}. Đặt
\begin{equation}
 q_\ell^{ab}=\mathbb E\!\left[z_i^{(\ell)}(\mathbf x_a)z_i^{(\ell)}(\mathbf x_b)\right],\qquad
 c_\ell=\frac{q_\ell^{ab}}{\sqrt{q_\ell^{aa}q_\ell^{bb}}}.
 \label{eq:long-corr}
\end{equation}
Giả sử phương sai của cả hai tín hiệu đã tiến gần cùng một điểm bất động $q_*>0$. Khi đó,
\begin{equation}
 \mathbb E\!\left[\big(z_i^{(\ell)}(\mathbf x_a)-z_i^{(\ell)}(\mathbf x_b)\big)^2\right]
 \approx2q_*(1-c_\ell).
 \label{eq:long-distance}
\end{equation}
Nếu $c_\ell\to1$, khoảng cách bình phương trung bình trên mỗi tọa độ của preactivation sẽ tiến về không, dù phương sai của từng tín hiệu vẫn xấp xỉ $q_*$. Với activation là Lipschitz như $\tanh$, khoảng cách giữa các activation cũng sẽ bị chặn theo khoảng cách của preactivation. Các biểu diễn cũng vì thế mà trở nên khó phân biệt hơn theo nghĩa khoảng cách này. Đây là sự suy giảm tín hiệu phân biệt đầu vào, không phải phát biểu rằng mọi activation đều tiến về không, cũng không phải một phép đo trực tiếp lượng thông tin Shannon.

\subsubsection{Liên hệ định lượng với sự suy giảm gradient}
Gần điểm bất động tương quan $c=1$, trong chế độ phương sai đã ổn định, tuyến tính hóa ánh xạ tương quan cho
\begin{equation}
 1-c_\ell\approx\chi(1-c_{\ell-1}),\qquad
 \chi=\sigma_w^2\int\mathcal Dz\,
 \left[\phi'\!\left(\sqrt{q_*}z\right)\right]^2.
 \label{eq:long-chi}
\end{equation}
Nếu $0<\chi<1$, sai khác giữa hai tín hiệu gần nhau giảm theo cấp số nhân quanh điểm bất động này. Độ sâu đặc trưng tương ứng là $\xi=-1/\log\chi$, bởi $\chi^d=e^{-d/\xi}$~\cite[\S2--3]{long-schoenholz2017}.

Đáng chú ý, cùng hệ số $\chi$ xuất hiện trong phân tích lan truyền ngược của tác giả. Với các lớp có cùng chiều rộng, đặt $\widetilde q_\ell=\mathbb E[(\delta_i^{(\ell)})^2]$. Dưới các xấp xỉ của mô hình,
\begin{equation}
 \widetilde q_\ell\approx\chi\widetilde q_{\ell+1},\qquad
 \widetilde q_\ell\approx\chi^{L-\ell}\widetilde q_L.
 \label{eq:long-mf-backward}
\end{equation}
Cụ thể, phân tích backward trong~\cite[\S4]{long-schoenholz2017} sử dụng thêm xấp xỉ độc lập giữa trọng số dùng trong forward và backward để khép kín quan hệ moment. Đây là công cụ phân tích, không phải quy trình backpropagation thực tế sử dụng hai bộ trọng số khác nhau. Khi các lớp có chiều rộng khác nhau, quan hệ truyền moment còn có hệ số tỷ lệ chiều rộng; khi mạng đã học, tương quan giữa trọng số và trạng thái cũng có thể làm xấp xỉ kém chính xác.

Trong phạm vi trên, $\chi<1$ gắn với sự suy giảm moment bậc hai của gradient, $\chi>1$ gắn với sự tăng trưởng, còn $\chi\approx1$ làm giảm biến dạng thang đo trung bình qua nhiều lớp. Tuy nhiên, ổn định một moment không bảo đảm mọi hướng gradient đều được bảo toàn, cũng không bảo đảm sẽ tối ưu hóa thành công. Kết quả này chủ yếu giải thích tại sao cấu hình khởi tạo ban đầu lại có thể đồng thời ảnh hưởng tới khả năng truyền tín hiệu xuôi và gradient ngược.

Từ góc độ đánh giá thực nghiệm, ta cần theo dõi đồng thời độ lớn hoặc moment bậc hai của activation, tỷ lệ neuron bão hòa, độ khác biệt giữa các biểu diễn và gradient trạng thái theo lớp, điều này giúp bổ sung một góc nhìn khác về \emph{hệ quả của phép biến đổi lên tín hiệu}; cách lựa chọn khởi tạo và độ sâu sẽ được trình bày ở các mục tương ứng của báo cáo.

\subsection{Vanishing Gradient trong RNN}
\label{sec:long-rnn}
\textit{Người phụ trách: Long}

Trong mạng nơ-ron hồi tiếp (\textit{recurrent neural network}, RNN), trạng thái ẩn được cập nhật nhiều lần bằng một hàm chuyển trạng thái dùng chung tham số. Vì vậy, ngay cả một RNN chỉ có một lớp hồi tiếp cũng tạo ra đồ thị tính toán sâu khi xử lý chuỗi dài. Vấn đề vanishing gradient xuất hiện theo \emph{khoảng cách thời gian}, làm suy yếu tín hiệu học liên quan đến những sự kiện ở xa thời điểm tính mất mát. Phân tích dưới đây dựa trên Pascanu và cộng sự~\cite[\S1.1--2.1]{long-pascanu2013} và Goodfellow và cộng sự~\cite[\S8.2.5, \S10.7]{long-goodfellow2016}.

\subsubsection{RNN và phép trải mạng theo thời gian}
Xét một RNN cơ bản có trạng thái ẩn kích thước $m$:
\begin{equation}
 \mathbf a_t=\mathbf W_h\mathbf h_{t-1}+\mathbf W_x\mathbf x_t+\mathbf b,
 \qquad \mathbf h_t=\phi(\mathbf a_t),\qquad
 \mathcal L=\sum_{t=1}^{T}\ell_t(\mathbf h_t).
 \label{eq:long-rnn-model}
\end{equation}
Hàm $\ell_t$ bao gồm phép đọc đầu ra và so sánh với nhãn tại thời điểm $t$. Trong phân tích này, trạng thái ban đầu $\mathbf h_0$ được giả sử cố định; các tham số hồi tiếp là $\boldsymbol\theta=\{\mathbf W_h,\mathbf W_x,\mathbf b\}$. Lan truyền ngược theo thời gian (\textit{backpropagation through time}, BPTT) thực hiện backpropagation trên đồ thị được trải thành $T$ bước. Các nút trạng thái thuộc những thời điểm khác nhau, nhưng sử dụng cùng một bộ tham số.

Jacobian của một bước chuyển là
\begin{equation}
 \mathbf A_t=\frac{\partial\mathbf h_t}{\partial\mathbf h_{t-1}}
 =\mathbf D_t\mathbf W_h,\qquad
 \mathbf D_t=\operatorname{diag}\!\left(\phi'(\mathbf a_t)\right).
 \label{eq:long-rnn-jac}
\end{equation}
Do $\mathbf a_t$ phụ thuộc vào đầu vào và trạng thái trước đó, các $\mathbf D_t$ thường khác nhau theo thời gian. Vì vậy, với RNN phi tuyến, tích Jacobian nói chung không thể thay bằng một lũy thừa của riêng $\mathbf W_h$.

\subsubsection{Tín hiệu sai số phải đi qua một tích Jacobian dài}
Với $t>k$, đặt
\begin{equation}
 \mathbf P_{t:k}=\frac{\partial\mathbf h_t}{\partial\mathbf h_k}
 =\mathbf A_t\mathbf A_{t-1}\cdots\mathbf A_{k+1},\qquad
 \mathbf P_{t:t}=\mathbf I.
 \label{eq:long-rnn-product}
\end{equation}
Nếu chỉ xét mất mát $\ell_t$ tại thời điểm $t$, đóng góp gradient của nó tại trạng thái $\mathbf h_k$ là
\begin{equation}
 \mathbf g_k^{(t)}
 :=\nabla_{\mathbf h_k}\ell_t
 =\mathbf P_{t:k}^{\mathsf T}\mathbf r_t,
 \qquad \mathbf r_t=\nabla_{\mathbf h_t}\ell_t.
 \label{eq:long-rnn-credit}
\end{equation}
Gradient từ sự kiện tại $t$ phải đi qua $t-k$ phép biến đổi để tác động đến trạng thái tại $k$. Với mất mát tại nhiều thời điểm, gradient tổng tại trạng thái thỏa mãn
\begin{equation}
 \mathbf g_k=\mathbf r_k+\mathbf A_{k+1}^{\mathsf T}\mathbf g_{k+1}
 =\sum_{t=k}^{T}\mathbf P_{t:k}^{\mathsf T}\mathbf r_t.
 \label{eq:long-rnn-total-state}
\end{equation}
Quan hệ này phân biệt tín hiệu sai số cục bộ với tín hiệu truyền từ các thời điểm xa. Một gradient tổng khác không chưa chứng minh rằng mạng vẫn nhận được đóng góp đáng kể từ các phụ thuộc dài hạn.

\subsubsection{Điều kiện đủ cho sự suy giảm theo cấp số nhân}
Giả sử tồn tại $\rho\in(0,1)$ sao cho $\|\mathbf A_s\|_2\leq\rho$ tại mọi bước trên đoạn đang xét. Tính dưới nhân của chuẩn phổ cho
\begin{equation}
 \|\mathbf g_k^{(t)}\|_2
 \leq\left(\prod_{s=k+1}^{t}\|\mathbf A_s\|_2\right)\|\mathbf r_t\|_2
 \leq\rho^{t-k}\|\mathbf r_t\|_2.
 \label{eq:long-rnn-bound}
\end{equation}
Khi gradient tại đầu cuối được giữ cố định hoặc bị chặn, ảnh hưởng truyền ngược suy giảm theo cấp số nhân của độ trễ. Nếu $|\phi'(u)|\leq\gamma$ với mọi $u$, thì
\begin{equation}
 \|\mathbf A_s\|_2\leq\gamma\|\mathbf W_h\|_2.
 \label{eq:long-rnn-sufficient}
\end{equation}
Do đó, $\gamma\|\mathbf W_h\|_2<1$ là một điều kiện đủ cho sự co đồng đều này. Với $\tanh$, có thể chọn $\gamma=1$; với sigmoid logistic, $\gamma=1/4$. Đối với $\tanh$, ngay cả khi $\|\mathbf W_h\|_2\geq1$, các trạng thái bão hòa vẫn có thể làm tích Jacobian co mạnh. Ngược lại, việc một số chuẩn Jacobian lớn hơn $1$ không đủ để kết luận gradient thực tế sẽ bùng nổ: kết quả còn phụ thuộc hướng của gradient và toàn bộ tích ma trận~\cite[\S2.1]{long-pascanu2013}.

Trường hợp tuyến tính giúp minh họa vai trò của phép lặp. Nếu $\phi(u)=u$, ta có $\mathbf P_{t:k}=\mathbf W_h^{t-k}$. Khi $\mathbf W_h$ chéo hóa được,
\begin{equation}
 \mathbf W_h=\mathbf V\boldsymbol\Lambda\mathbf V^{-1}
 \quad\Longrightarrow\quad
 \mathbf W_h^{t-k}=\mathbf V\boldsymbol\Lambda^{t-k}\mathbf V^{-1}.
 \label{eq:long-rnn-eigen}
\end{equation}
Các thành phần ứng với $|\lambda_i|<1$ suy giảm khi độ trễ tăng. Tổng quát hơn, bán kính phổ $r(\mathbf W_h)<1$ bảo đảm $\mathbf W_h^d\to\mathbf0$ khi $d\to\infty$, kể cả khi không chéo hóa được. Tuy nhiên, một ma trận không chuẩn tắc vẫn có thể gây khuếch đại tạm thời trước khi suy giảm tiệm cận. Với RNN phi tuyến, cần quay lại tích $\mathbf A_t\cdots\mathbf A_{k+1}$; chỉ xét eigenvalue của ma trận trọng số là chưa đủ.

\subsubsection{Gradient tham số dùng chung và bài toán gán trách nhiệm}
Để hiểu tác động lên việc học, viết $\boldsymbol\theta$ thành vector các tham số hồi tiếp và đặt
\begin{equation}
 \mathbf B_k=\left.\frac{\partial\mathbf h_k}{\partial\boldsymbol\theta}
 \right|_{\mathbf h_{k-1}\ \text{cố định}}.
 \label{eq:long-rnn-local}
\end{equation}
$\mathbf B_k$ chỉ ghi nhận tác động trực tiếp của tham số tại bước $k$. Vì các tham số được dùng lại ở mọi thời điểm, quy tắc chuỗi cho
\begin{equation}
 \nabla_{\boldsymbol\theta}\mathcal L
 =\sum_{t=1}^{T}\sum_{k=1}^{t}
 \mathbf B_k^{\mathsf T}\mathbf P_{t:k}^{\mathsf T}\mathbf r_t.
 \label{eq:long-rnn-param}
\end{equation}
Đây là dạng tổng các đóng góp theo thời gian trong phân tích của Pascanu và cộng sự~\cite[\S1.1]{long-pascanu2013}. Mỗi số hạng biểu diễn tác động của việc sử dụng tham số tại bước $k$ đến mất mát ở bước $t$. Nếu $\|\mathbf B_k\|_2\leq M$ và điều kiện của~\eqref{eq:long-rnn-bound} được thỏa mãn thì
\begin{equation}
 \left\|\mathbf B_k^{\mathsf T}\mathbf P_{t:k}^{\mathsf T}\mathbf r_t\right\|_2
 \leq M\rho^{t-k}\|\mathbf r_t\|_2.
 \label{eq:long-rnn-term}
\end{equation}
Các đóng góp có độ trễ lớn vì thế có thể trở nên rất mờ nhạt trong khi các đóng góp gần thời điểm đầu ra vẫn đáng kể. Điều này giải thích vì sao mạng vẫn có thể giảm loss và học các quan hệ ngắn hạn, nhưng lại khó gán trách nhiệm cho những sự kiện ở thời điểm xa hơn. Vanishing gradient trong RNN không nhất thiết có nghĩa là toàn bộ $\nabla_{\boldsymbol\theta}\mathcal L$ gần bằng không.

Về mặt lan truyền xuôi, nếu hai chuỗi có lịch sử khác nhau tại bước $k$ nhưng lại cùng các đầu vào từ $k+1$ đến $t$, thì trong xấp xỉ cục bộ,
ta vẫn có $\Delta\mathbf h_t\approx\mathbf P_{t:k}\Delta\mathbf h_k$.
Do đó, cùng tích Jacobian có thể làm giảm cả khác biệt do lịch sử để lại và tín hiệu sai số quay về lịch sử đó. Tuy vậy, năng lực biểu diễn hoặc lưu giữ thông tin và khả năng học cách sử dụng thông tin bằng gradient là hai vấn đề khác nhau; không nên đồng nhất vanishing gradient với một khẳng định tuyệt đối rằng RNN không thể ghi nhớ xa~\cite[\S10.7]{long-goodfellow2016}.

\subsubsection{Ví dụ định lượng và hệ quả khi đánh giá mô hình}
Xét RNN một chiều $h_s=\tanh(wh_{s-1}+ux_s+b)$. Đạo hàm qua $d$ bước là
\begin{equation}
 \frac{\partial h_t}{\partial h_{t-d}}
 =\prod_{s=t-d+1}^{t}w(1-h_s^2).
 \label{eq:long-rnn-scalar}
\end{equation}
Chọn $w=0.9$, $b=0$, trạng thái ban đầu bằng không và đầu vào bằng không trên đoạn xét. Khi đó $h_s=0$ ở mọi bước, nên đạo hàm trên bằng chính xác $0.9^d$. Bảng~\ref{tab:long-decay} cho thấy sự suy giảm ngay cả khi $\tanh$ hoạt động tại vùng có đạo hàm lớn nhất.
\begin{table}[htbp]
 \centering
 \caption{Hệ số truyền độ nhạy theo độ trễ trong RNN một chiều. Các giá trị được tính trực tiếp từ $0.9^d$, không phải kết quả huấn luyện.}
 \label{tab:long-decay}
 \begin{tabular}{@{}rrrrr@{}}
 \toprule
 Độ trễ $d$ & $10$ & $20$ & $50$ & $100$ \\
 \midrule
 $0.9^d$ & $0.3487$ & $0.1216$ & $5.154\times10^{-3}$ & $2.656\times10^{-5}$ \\
 \bottomrule
 \end{tabular}
\end{table}
Nếu các trạng thái đi vào vùng bão hòa, những hệ số $1-h_s^2$ nhỏ còn làm độ nhạy suy giảm nhanh hơn. Ví dụ này minh họa sự suy giảm gradient trạng thái; nó không chứng minh mọi gradient tham số đều khác không hoặc đều suy giảm cùng tốc độ.

Để quan sát đúng cơ chế trong thí nghiệm, có thể đặt mất mát tại bước cuối và đo $\|\nabla_{\mathbf h_k}\ell_T\|_2$ theo độ trễ $T-k$, đồng thời đánh giá hiệu năng trên một tác vụ cần thông tin từ đầu chuỗi. Nếu đồ thị tính toán bị cắt tại biên của \textit{truncated BPTT}, đóng góp vượt ngoài cửa sổ bị loại bỏ do cách tính, cần phân biệt với gradient suy giảm qua một đồ thị còn nguyên. Tương tự, gradient clipping giới hạn gradient quá lớn nhưng không khôi phục những đóng góp dài hạn đã mất~\cite[\S10.11.1]{long-goodfellow2016}. Các kiến trúc có đường truyền trạng thái được điều tiết bằng cổng, như LSTM và GRU, sẽ được bàn trong mục giải pháp; ở đây, kết luận cốt lõi là chiều dài chuỗi làm tăng số Jacobian phải nhân và có thể khiến tín hiệu gán trách nhiệm dài hạn suy yếu nghiêm trọng.

\subsection{Ảnh hưởng của kiến trúc mạng}

Kiến trúc mạng quyết định cách tín hiệu được truyền đi và những đường mà
gradient có thể đi qua. Đối với mạng tuần tự, đầu ra của mỗi lớp chỉ được
truyền trực tiếp tới lớp kế tiếp. Do đó, gradient từ hàm mất mát muốn lan
truyền ngược về các lớp gần đầu vào phải đi qua toàn bộ các phép biến đổi nằm
giữa hai lớp. Theo chain rule, đường truyền này chứa tích của nhiều ma trận
Jacobian. Khi nhiều Jacobian trong tích có chuẩn nhỏ hơn $1$, gradient sẽ suy
giảm nhanh trước khi đến được các lớp đầu tiên. Vì vậy, ngoài hàm activation
và cách khởi tạo trọng số như đã trình bày ở trên, độ sâu, độ rộng và cách kết
nối giữa các lớp cũng ảnh hưởng trực tiếp đến khả năng tối ưu của mô hình.


\subsubsection{Tương quan giữa độ sâu và độ rộng}

Theo paper của Hanin (2018)~\cite{hanin2018}  mạng fully connection sử dụng activatoin ReLU.
Trọng số giữa các lớp được giả sử là độc lập, đối xứng quanh 0 và có phương sai tỉ lệ phù hợp với
ReLU. Trong phạm vi mô hình này, tác giả đã chỉ ra rằng ảnh hưởng kết hợp của độ sâu và độ rộng tác động
đến sự phân tán của gradient được đạc trưng bởi đại lượng $\beta$.
\[
    \beta=\sum_{j=1}^{d-1}\frac{1}{n_j},
\]

\begin{infobox}{Ý nghĩa của đại lượng $\beta$}
\small
\begin{description}[
    style=multiline,
    leftmargin=3.3em,
    labelwidth=2.5em,
    labelsep=0.8em,
    itemsep=3pt,
    topsep=2pt
]
    \item[$\beta$] Là một hằng số phụ thuộc vào kiến trúc, dùng để tóm tắt ảnh hưởng
    của sự kết hợp về độ sâu và độ rộng lên mức độ phân tán của các phần tử trong ma trận
    Jacobian ở thời điểm khởi tạo.
    \item[$d$] số hidden layer của mô hình.

    \item[$j$] Chỉ số của lớp ẩn.

    \item[$n_j$] Số neuron, hay độ rộng, của lớp ẩn thứ $j$.

    \item[$1/n_j$] Mức đóng góp của lớp thứ $j$ vào $\beta$. Lớp càng hẹp
    thì $1/n_j$ càng lớn; lớp càng rộng thì đóng góp này càng nhỏ.
\end{description}

Trong phạm vi mạng fully connected dùng ReLU được paper phân tích,
$\beta$ nhỏ tương ứng với các phần tử Jacobian ổn định hơn về độ phân tán;
$\beta$ lớn cho thấy nguy cơ xuất hiện đồng thời nhiều đạo hàm rất nhỏ và rất
lớn tại thời điểm khởi tạo.
\end{infobox}

Để nghiên cứu quá trình truyền gradient, paper xét phần tử của Jacobian đầu
vào--đầu ra
\[
    Z_{p,q}=\frac{\partial f_q(\mathbf{x})}{\partial x_p},
\]

\begin{infobox}{Ý nghĩa các ký hiệu trong $Z_{p,q}$}
\small
\begin{description}[
    style=multiline,
    leftmargin=3.8em,
    labelwidth=3em,
    labelsep=0.8em,
    itemsep=3pt,
    topsep=2pt
]
    \item[$\mathbf{x}$] Vector đầu vào của mạng.

    \item[$x_p$] Phần tử thứ $p$ của vector đầu vào.

    \item[$f_q(\mathbf{x})$] Đầu ra thứ $q$ mà mạng tính được từ
    $\mathbf{x}$.

    \item[$p,q$] Chỉ số dùng để chọn một tọa độ đầu vào và một tọa độ đầu ra.

    \item[$Z_{p,q}$] Một phần tử của Jacobian đầu vào--đầu ra, đo độ nhạy của
    đầu ra thứ $q$ khi đầu vào thứ $p$ thay đổi một lượng rất nhỏ. Nếu
    $|Z_{p,q}|$ gần $0$, ảnh hưởng theo cặp đầu vào--đầu ra này gần như biến
    mất; nếu nó quá lớn, ảnh hưởng bị khuếch đại mạnh.
\end{description}
\end{infobox}

Với cách định tỉ lệ phương sai nói trên, moment bậc hai của phần tử Jacobian
được giữ ổn định:
\[
    \mathbb{E}\!\left[Z_{p,q}^{2}\right]=\frac{1}{n_0},
\]

Trong biểu thức này, $\mathbb{E}[\,\cdot\,]$ là kỳ vọng qua các lần khởi tạo
ngẫu nhiên, còn $n_0$ là số chiều của vector đầu vào. Phương trình cho biết
bình phương độ nhạy được giữ ổn định về trung bình; nó không khẳng định mọi
$Z_{p,q}$ đều có độ lớn gần nhau. Hanin chứng minh
rằng moment bậc cao và độ phân tán của bình phương các phần tử Jacobian tăng
theo dạng hàm mũ của $\beta$. Khi $\beta$ lớn, trong cùng một mạng có thể tồn
tại đồng thời các đạo hàm rất nhỏ và các đạo hàm rất lớn, khiến một learning
rate chung khó phù hợp với tất cả tham số.

Nếu mọi lớp ẩn có cùng độ rộng $n$, đại lượng trên trở thành
\[
    \beta=\frac{d-1}{n}.
\]
Công thức này cho thấy độ sâu không nên được xem xét độc lập với độ rộng. Nếu
giữ nguyên $n$ nhưng tiếp tục tăng số lớp, $\beta$ tăng và gradient tại thời
điểm khởi tạo có xu hướng phân tán mạnh hơn. Ngược lại, tăng độ rộng tương ứng
với độ sâu giúp hạn chế sự tích lũy của các hiệu ứng hữu hạn chiều rộng.

Trong thực nghiệm, xu hướng tăng độ rộng theo chiều sâu thường xuất hiện rõ ở
các mạng CNN. Mạng thường được chia thành nhiều
\textit{stage}; mỗi khi kích thước không gian của feature map giảm, số kênh
được tăng lên, phổ biến theo dãy $64\rightarrow128\rightarrow256\rightarrow512$.
Chẳng hạn, ResNet tăng gấp đôi số kênh tại các điểm chuyển stage
\cite{he2016identity}. Cách thiết kế này trước hết giúp các lớp sâu có đủ khả
năng biểu diễn và giữ chi phí tính toán tương đối cân bằng khi chiều cao và
chiều rộng của feature map giảm. Nếu liên hệ với trực giác từ $\beta$, việc
không để các lớp ẩn trở nên quá hẹp cũng làm giảm đóng góp $1/n_j$ của từng
lớp. Tuy nhiên, quy tắc ``gấp đôi số kênh'' là một kinh nghiệm thiết kế CNN,
không phải hệ quả trực tiếp từ công thức $\beta$ nhưng nó cũng phần chứng mình
được quá trình thực nghiệm cũng chỉ ra được xu hướng cũng như là điểm nghẽn này.

Ở gần đầu ra, độ rộng thường giảm mạnh thông qua global average pooling hoặc
một lớp tuyến tính ánh xạ đặc trưng sang số lớp cần dự đoán. Điều này không mâu
thuẫn với phân tích trên, vì đây chủ yếu là bước đọc kết quả cuối cùng chứ không
phải một chuỗi dài các lớp ẩn hẹp. Hơn nữa, số chiều đầu ra không xuất hiện
trong tổng $\beta=\sum_{j=1}^{d-1}1/n_j$. Trường hợp đáng lưu ý hơn là một
bottleneck rất hẹp nằm giữa đường truyền sâu, vì nó trực tiếp làm tăng $\beta$
và có thể trở thành điểm nghẽn đối với sự lan truyền gradient.

\subsubsection{Ảnh hưởng của lớp bottleneck}

Do mỗi lớp đóng góp một lượng $1/n_j$ vào $\beta$, một lớp quá hẹp có thể chi
phối đáng kể đại lượng này. Chẳng hạn, xét hai mạng có hai lớp ẩn và cùng tổng
số neuron ẩn:
\[
    \beta_{(50,50)}
    =\frac{1}{50}+\frac{1}{50}=0.04,
    \qquad
    \beta_{(10,90)}
    =\frac{1}{10}+\frac{1}{90}\approx0.111.
\]
Mặc dù cả hai kiến trúc đều có tổng cộng $100$ neuron ẩn, kiến trúc thứ hai có
$\beta$ lớn hơn do chứa một lớp hẹp. Với độ sâu và tổng số neuron ẩn cố định,
phân bố độ rộng tương đối đồng đều giữa các lớp sẽ làm giảm $\beta$. Tuy nhiên,
kết quả này không có nghĩa rằng mọi bottleneck đều gây vanishing gradient; nó
chỉ mô tả sự phân tán của các phần tử Jacobian trong mạng fully connected dùng
ReLU tại thời điểm khởi tạo, dưới các giả thiết của paper.

\subsubsection{Ảnh hưởng của cách kết nối giữa các lớp}

Độ sâu và độ rộng chưa mô tả đầy đủ kiến trúc mạng. Trong plain network,
gradient chỉ có đường đi tuần tự qua các lớp trung gian. Ngược lại, skip
connection tạo thêm đường truyền ngắn hơn, cho phép tín hiệu và gradient bỏ qua
một hoặc nhiều phép biến đổi. Đối với residual block có dạng
\[
    \mathbf{h}^{(l+1)}
    =\mathbf{h}^{(l)}
    +\mathcal{F}\!\left(\mathbf{h}^{(l)},\mathbf{W}^{(l)}\right),
\]
Jacobian của block theo đầu vào là
\[
    \frac{\partial\mathbf{h}^{(l+1)}}{\partial\mathbf{h}^{(l)}}
    =\mathbf{I}
    +\frac{\partial\mathcal{F}}{\partial\mathbf{h}^{(l)}}.
\]
Thành phần đồng nhất $\mathbf{I}$ tạo một đường truyền trực tiếp, nhờ đó
gradient không hoàn toàn phụ thuộc vào tích các đạo hàm trên nhánh biến đổi
$\mathcal{F}$. Kết quả phân tích của He và cộng sự
(2016)~\cite{he2016identity} cũng cho thấy đường truyền đồng nhất đóng vai trò
quan trọng đối với quá trình lan truyền xuôi và lan truyền ngược. Dù vậy, hiệu
quả còn phụ thuộc vào các phép toán đặt trên nhánh
tắt; một shortcut có phép co giãn, cổng điều khiển hoặc phép chiếu vẫn có thể
cản trở sự truyền tín hiệu. Cơ chế và cách sử dụng residual/skip connection sẽ
được trình bày chi tiết hơn trong phần giải pháp.

\begin{notebox}
Kiến trúc ảnh hưởng đến gradient thông qua độ dài đường truyền, tương quan
giữa độ sâu và độ rộng, và cách các lớp được kết nối. Trong mạng fully
connected dùng ReLU tại khởi tạo, $\beta=\sum_j 1/n_j$ cung cấp một đại lượng
định lượng cho ảnh hưởng kết hợp của độ sâu và độ rộng; với residual network,
identity shortcut còn tạo thêm một thành phần truyền gradient trực tiếp.
\end{notebox}

\section{Ví dụ minh họa cho các trường hợp Vanishing Gradient}
\textit{Người phụ trách chung: Khanh, Tâm, Long, Khoa}

\subsection{Khung mô hình mô-đun cho thí nghiệm ablation}

Để phân tích riêng tác động của từng thành phần, có thể xây dựng một backbone
theo dạng mô-đun như Hình~\ref{fig:modular-vanishing-ablation}. Sáu nhóm biến
được đánh số từ 1 đến 6 có thể thay đổi độc lập, gồm độ sâu, độ rộng hoặc
bottleneck, activation, phương pháp khởi tạo trọng số, normalization và
residual connection. Các vị trí mang cùng số trong backbone cho biết thành
phần nào của mô hình chịu ảnh hưởng bởi biến tương ứng.

\begin{figure}[H]
    \centering
    \includegraphics[width=\linewidth]{assets/modular_vanishing_gradient_ablation.pdf}
    \caption{Mô hình mô-đun dùng để tách riêng ảnh hưởng của các thành phần đến
    vanishing gradient. Các hộp màu cam là biến được can thiệp; đường màu xanh
    dương là đường truyền tín hiệu; các đầu dò màu xanh lá ghi nhận gradient
    tại nhiều độ sâu. Sơ đồ được thiết kế bằng Draw.io.}
    \label{fig:modular-vanishing-ablation}
\end{figure}

Mỗi thí nghiệm bắt đầu từ cùng một cấu hình baseline và chỉ thay đổi một hộp
màu cam, trong khi giữ cố định dữ liệu, optimizer, learning rate, batch size,
số epoch và tập seed. Tại mỗi lớp, cần ghi lại gradient norm, histogram
gradient hoặc Jacobian, đồng thời theo dõi loss và accuracy. So sánh các đường
cong này giữa baseline và cấu hình đã thay đổi sẽ giúp quy tác động quan sát
được cho đúng thành phần, thay vì thay nhiều yếu tố cùng lúc rồi không xác định
được nguyên nhân.

\subsubsection{Bộ dữ liệu CIFAR-10}

Các thí nghiệm trong phần này sử dụng CIFAR-10, bộ dữ liệu phân loại ảnh được
giới thiệu trong báo cáo của Krizhevsky (2009)~\cite{krizhevsky2009}. CIFAR-10
gồm $60\,000$ ảnh màu RGB kích thước $32\times32$, được chia đều vào 10 lớp:
\textit{airplane}, \textit{automobile}, \textit{bird}, \textit{cat},
\textit{deer}, \textit{dog}, \textit{frog}, \textit{horse}, \textit{ship} và
\textit{truck}. Mỗi lớp có $6\,000$ ảnh. Phân chia chính thức gồm $50\,000$ ảnh
huấn luyện và $10\,000$ ảnh kiểm tra.

Trong báo cáo này, $5\,000$ ảnh được tách có phân tầng từ tập huấn luyện chính
thức để làm tập validation, tức mỗi lớp đóng góp $500$ ảnh. Vì vậy, dữ liệu
được sử dụng theo tỉ lệ $45\,000/5\,000/10\,000$ cho train, validation và test;
tập test chỉ được dùng khi đánh giá mô hình cuối cùng. Ảnh đầu vào được chuyển
sang số thực, chuẩn hóa riêng trên từng kênh RGB bằng mean và standard deviation
tính từ tập train. Random crop với padding 4 và random horizontal flip chỉ được
áp dụng cho tập train. Kích thước ảnh nhỏ và số lượng mẫu vừa phải giúp có thể
lặp lại nhiều cấu hình độ sâu, độ rộng và skip connection với cùng một quy
trình huấn luyện.

\subsubsection{Cấu hình mô hình baseline}

Baseline được chọn làm điểm tham chiếu ổn định cho Hình~\ref{fig:modular-vanishing-ablation}.
Đầu tiên, một convolution $3\times3$ ánh xạ ảnh RGB sang $64$ kênh. Backbone
gồm $L=8$ block có cùng độ rộng $n_l=64$. Trong mỗi block, nhánh residual thực
hiện lần lượt convolution $3\times3$, Batch Normalization và ReLU, sau đó cộng
với identity shortcut. Cuối mạng, global average pooling tạo một vector $64$
chiều và lớp tuyến tính ánh xạ vector này sang 10 logit. Cấu hình đầy đủ được
khai báo trong Bảng~\ref{tab:cifar10-baseline}.

\begin{table}[H]
    \centering
    \caption{Cấu hình baseline cho thí nghiệm trên CIFAR-10.}
    \label{tab:cifar10-baseline}
    \small
    \renewcommand{\arraystretch}{1.18}
    \begin{tabularx}{\linewidth}{
        >{\raggedright\arraybackslash\bfseries}p{0.23\linewidth}
        >{\raggedright\arraybackslash}X}
        \toprule
        Thành phần & Giá trị baseline \\
        \midrule
        Dữ liệu
            & CIFAR-10; $45\,000$ train, $5\,000$ validation và $10\,000$ test. \\
        Đầu vào
            & Tensor $3\times32\times32$; chuẩn hóa theo từng kênh RGB. \\
        Data augmentation
            & Random crop $32\times32$ với padding 4 và random horizontal flip
              cho tập train. \\
        Stem
            & Convolution $3\times3$, stride 1, padding 1, từ 3 sang 64 kênh. \\
        Độ sâu và độ rộng
            & $L=8$ block; mọi block có $n_l=64$ kênh; không có bottleneck. \\
        Residual block
            & $\operatorname{Conv}_{3\times3}\rightarrow\operatorname{BatchNorm}
              \rightarrow\operatorname{ReLU}\rightarrow+$ identity shortcut. \\
        Khởi tạo
            & He normal cho convolution và lớp tuyến tính; bias bằng 0. \\
        Đầu ra và loss
            & Global average pooling, Linear $64\rightarrow10$ và cross-entropy loss. \\
        Optimizer
            & SGD, learning rate $0.1$, momentum $0.9$, weight decay $5\times10^{-4}$. \\
        Lịch huấn luyện
            & 100 epoch, batch size 128 và cosine learning-rate schedule. \\
        Số lần lặp
            & Ba seed độc lập: $0$, $1$ và $2$; báo cáo mean $\pm$ standard deviation. \\
        Đại lượng theo dõi
            & $\lVert\partial\mathcal{L}/\partial\mathbf{h}^{(l)}\rVert_2$,
              $\lVert\partial\mathcal{L}/\partial\mathbf{W}^{(l)}\rVert_2$,
              tỉ lệ gradient lớp đầu/lớp cuối, loss và accuracy. \\
        \bottomrule
    \end{tabularx}
\end{table}

Từ baseline trên, mỗi nhóm thí nghiệm chỉ thay đổi một thành phần được đánh số
trong Hình~\ref{fig:modular-vanishing-ablation}:
\begin{enumerate}[label=\textbf{\arabic*.}, leftmargin=2.2em, itemsep=3pt]
    \item \textbf{Độ sâu:} $L\in\{4,8,16,32\}$, giữ độ rộng bằng 64.
    \item \textbf{Độ rộng:} $n_l\in\{32,64,128\}$; bổ sung một cấu hình có
    bottleneck 16 kênh ở giữa mạng.
    \item \textbf{Activation:} Sigmoid, Tanh hoặc ReLU.
    \item \textbf{Khởi tạo:} Gaussian
    $\mathcal{N}(0,0.01^2)$, Xavier hoặc He initialization.
    \item \textbf{Normalization:} không dùng normalization, BatchNorm hoặc
    LayerNorm.
    \item \textbf{Shortcut:} tắt residual connection, dùng identity shortcut,
    hoặc dùng projection shortcut khi kích thước hai nhánh khác nhau.
\end{enumerate}

Khi thay đổi độ sâu hoặc độ rộng, số tham số và chi phí tính toán cũng thay đổi;
do đó hai đại lượng này cần được báo cáo cùng kết quả. Với các nhóm còn lại,
kiến trúc và số tham số được giữ nguyên. Dữ liệu, thứ tự mini-batch, optimizer,
learning-rate schedule và tập seed phải giống baseline để khác biệt trong
gradient có thể được quy về thành phần đang khảo sát.

\subsection{MLP sâu dùng sigmoid}
Xét mạng Perceptron đa lớp (MLP) gồm $L$ lớp sử dụng hàm kích hoạt Sigmoid
\begin{equation}
    \sigma(z)=\frac{1}{1+e^{-z}}.
\end{equation}
Đạo hàm của hàm mất mát $\mathcal{L}$ đối với ma trận trọng số $W\_s$ tại lớp thứ $s$ được xác định theo công thức lan truyền ngược:
\begin{equation}
    \frac{\partial \mathcal{L}}{\partial W_s} = \left[ \text{diag}(d_s) B_s(E) \right] \otimes x_s^T, \quad s= 0, 1, \dots, L-1
\end{equation}
Trong đó $E = \frac{\partial \mathcal{L}}{\partial x_L}$ là vector lỗi tại đầu ra, $d_s= \sigma'(y_s)$ là vector đạo hàm kích hoạt tại lớp $s$, và toán tử truyền ngược tín hiệu lỗi $B_s(E)$ được tính thông qua ma trận Jacobian:
\begin{equation}
    B_s(E)= \left( \frac{\partial x_l}{\partial x_{s+1}} \right)^T E = \left(\prod_{l = s+1}^{L-1} \text{diag}(d_l) W_l \right)^T E
\end{equation}

Khi chiều sâu $L$ của mạng tăng lên, biến động về chuẩn đạo hàm (gradient norm) giữa lớp đầu tiên ($ s= 0$) và lớp cuối cùng ($s = L - 1$) thể hiện sự bất đối xứng rõ rệt:

\begin{enumerate}
    \item \textbf{Gradient Norm tại Lớp Cuối ($s = L -1$):} Tại lớp kề sát đầu ra, tín hiệu lỗi không phải truyền qua bất kỳ ma trận Jacobian trung gian nào, do đó $B_{L-1}(E) = E$. Chuẩn đạo hàm của lớp cuối phụ thuộc trực tiếp vào lỗi đầu ra $E$ và vector kích hoạt $x_{L-1}$:
    \begin{equation}
        \left\| \frac{\partial \mathcal{L}}{\partial W_{L-1}} \right\|_F = \left\| E \right\|_2 = \left\| \text{diag}(d_{L-1}) B_{L-1}(E) \right\|_2.
    \end{equation}
    Biên độ này duy trì ổn định ở mức hữu hạn và hoàn toàn độc lập với chiều sâu tổng thể $L$ của mạng.
    \item \textbf{Gradient Norm tại Lớp Đầu Tiên ($s = 0$):} Tín hiệu lỗi $E$ khi truyền ngược về lớp đầu tiên phải nhân nối tiếp qua chuỗi $L-1$ ma trận Jacobian:
    \begin{equation}
        \frac{\partial x_L}{\partial x_1} = \prod_{l = 1}^{l - 1} \text{diag}(d_l) W_l
    \end{equation}
    Do đạo hàm Sigmoid bị chặn trên bởi $r = \max_y \sigma'(y) = 0.25$, ta có $\|\text{diag}(d_l)\|_2 \le 0.25$. Đặt $\Omega_{\max} = \max_l \|W_l\|2$, chuẩn của ma trận Jacobian toàn cục thỏa mãn bất đẳng thức:
    \begin{equation}
        \left\| \frac{\partial x_L}{\partial x_1} \right\|_2 \le \prod_{l=1}^{L-1} \|\text{diag}(d_l)\|_2 \|W_l\|_2 \le (0.25 \cdot \Omega_{\max})^{L-1}
    \end{equation}
    Khi $0.25 \cdot \Omega_{\max} < 1$, chuẩn ma trận Jacobian suy giảm $0$ theo tốc độ hàm mũ khi $L \to \infty$ Do đó, gradient norm ở lớp đầu bị triệt tiêu tiệm cận:
    \begin{equation}
        \lim_{K \to \infty} \left\| \frac{\partial \mathcal{L}}{\partial W_0} \right\|_F = 0
    \end{equation}

    \item \textbf{So sánh Tỷ lệ Gradient Norm theo Chiều sâu $L$:}
    \begin{itemize}
        \item \textbf{Mạng nông ($L$ nhỏ):} Tỷ lệ giữa gradient norm lớp đầu và lớp cuối xấp xỉ tương đương, đảm bảo tất cả các lớp được cập nhật đồng đều.
        \item \textbf{Mạng sâu ($L \to \infty$):} Tỷ lệ gradient norm giữa lớp đầu và lớp cuối suy giảm theo hàm số mỹ về $0$:
        \begin{equation}
            \lim_{L \to \infty} \frac{\left\| \frac{\partial \mathcal{L}}{\partial W_0} \right \|_F}{\left | \frac{\partial \mathcal{L}}{\partial W_{L-1}} \right\|F} = 0
        \end{equation}
        Sự phân hóa này làm cho trộng số ở lớp cuối vẫn liên tục được điều chỉnh, trong khi các trọng số ở lớp đầu tiên rơi vào trạng thái đóng băng vĩnh viễn (gradient immobility), khiến mạng không thể học các biểu diển đặc trưng ở các lớp ban đầu.
    \end{itemize}
\end{enumerate}

\subsection{Ảnh hưởng của khởi tạo weight và độ sâu}
\textit{Phụ trách: Khoa}

Giữ nguyên dữ liệu và kiến trúc cơ bản, thay đổi độ sâu cùng phương pháp khởi tạo để quan sát loss và gradient.

\subsection{RNN học phụ thuộc dài hạn}
\textit{Phụ trách: Long}

Dùng một bài toán chuỗi đơn giản, thay đổi sequence length và đo gradient theo khoảng cách thời gian.

\subsection{Mạng có và không có residual connection}
\textit{Phụ trách: Tâm}

So sánh đường truyền gradient của plain network với residual network có cùng độ sâu.

\begin{examplebox}{Đầu ra nên có cho mỗi ví dụ}
Mỗi ví dụ nên gồm: mục tiêu, kiến trúc, cấu hình huấn luyện, đoạn code tối giản, biểu đồ gradient/loss và phần nhận xét ngắn gọn.
\end{examplebox}

\section{Các giải pháp phổ biến}
\textit{Người phụ trách chung: Khanh, Tâm, Long, Khoa}

\subsection{Lựa chọn hàm activation phù hợp}
\textit{Phụ trách: Khanh}

ReLU, Leaky ReLU, ELU/GELU và lưu ý về \textit{dying ReLU}.

\subsection{Khởi tạo trọng số và chuẩn hóa}
\label{subsec:fix-weightinit-normalization}
% \textit{Phụ trách: Khoa}

% Xavier/Glorot, He initialization, Batch Normalization và Layer Normalization.

Mục V.2 đã chỉ ra rằng gradient sẽ bị triệt tiêu khi hệ số khuếch đại $\chi$ nhỏ hơn 1. Các quy tắc Xavier/Glorot và He lựa chọn $\operatorname{Var}(W)$ sao cho $\chi \approx 1$.

\subsubsection{Khởi tạo Xavier/Glorot}
Theo Glorot và Bengio (2010)~\cite{glorot2010}, đối với một hàm kích hoạt gần như tuyến tính quanh $0$, phương sai của tín hiệu lan truyền thuận và của sai số lan truyền ngược thay đổi qua từng lớp theo công thức:
\begin{align*}
    \operatorname{Var}\big( \mathbf{z}^{(k)} \big)
& = n_{\text{in}} \operatorname{Var}(W)\, \operatorname{Var}\big( \mathbf{z}^{(k-1)} \big), \\
\operatorname{Var}\big( \boldsymbol\delta^{(k-1)} \big)&  = n_{\text{out}} \operatorname{Var}(W)\, \operatorname{Var}\big( \boldsymbol \delta^{(k)} \big).
\end{align*}
Để giữ phương sai lan truyền thuận không đổi, ta cần $n_{\text{in}} \operatorname{Var}(W) = 1$, và để giữ phương sai lan truyền ngược không đổi, ta cần $n_{\text{out}} \operatorname{Var}(W) = 1$. Hai điều kiện chỉ đồng thời thỏa mãn khi $n_{\text{in}} =n_{\text{out}}$, vì vậy bài báo sử dụng một lựa chọn dung hòa
\[ n_{\text{avg}} = \frac{n_{\text{in}} + n_{\text{out}}}{2}\]
Từ đây ta có \textit{khởi tạo Xavier/Glorot}:
\begin{itemize}
    \item Trọng số được khởi tạo từ một phân phối chuẩn với kỳ vọng 0 và phương sai $\sigma^2 = \dfrac{1}{n_{\text{avg}}}$.

    \item Trọng số được khởi tạo từ một phân phối đều trên đoạn từ $-r$ đến $r$, với $r = \sqrt{\dfrac{3}{n_{\text{avg}}}}$
\end{itemize}
và được sử dụng khi hàm kích hoạt là sigmoid hoặc tanh.

\subsubsection{Khởi tạo He/Kaiming}

Theo He và cộng sự (2015)~\cite{he2015rectifiers}, hàm ReLU đưa một nửa số giá trị tiền kích hoạt (pre-activation) về 0, và điều này bổ sung một hệ số $\frac{1}{2}$ vào công thức lan truyền tiến:
\[
\operatorname{Var}\big( \mathbf{z}^{(k)} \big) = \tfrac{1}{2}\, n_{\text{in}} \operatorname{Var}(W)\, \operatorname{Var}\big( \mathbf{z}^{(k-1)} \big)
\quad\Longrightarrow\quad
\operatorname{Var}(W) = \frac{2}{n_{\text{in}}},
\]
Từ đây ta có \textit{khởi tạo He/Kaiming:}
\begin{itemize}
    \item Trọng số được khởi tạo từ một phân phối chuẩn với kỳ vọng 0 và phương sai $\sigma^2 = \dfrac{2}{n_{\text{in}}}$.
\end{itemize}
và được sử dụng cho hàm ReLU và các biến thể của nó (Leaky ReLU, ELU, ...). He và cộng sự (2015)~\cite{he2015rectifiers} đã báo cáo rằng quy tắc này giúp huấn luyện từ đầu các mô hình kích hoạt ReLU cực kỳ sâu.

Với một trong hai quy tắc Xavier và He, độ lớn của các giá trị kích hoạt và gradient không bị thu nhỏ cũng như không bị phóng đại qua từng lớp tại thời điểm khởi tạo. Tích của các ma trận Jacobian tại Mục IV.1 khi đó không còn giảm dần theo số lượng lớp, ít nhất là tại thời điểm bắt đầu huấn luyện.

Tuy nhiên, hai quy tắc này chỉ khắc phục điểm khởi đầu vì chúng giả định các lớp có chiều rộng lớn và trọng số độc lập, và chỉ mô tả hành vi trung bình, và các trọng số sẽ thay đổi khi quá trình huấn luyện bắt đầu. Trong khi đó, chuẩn hóa (normalization) phát huy tác dụng xuyên suốt quá trình huấn luyện.

\subsubsection{Batch Normalization (Chuẩn hóa theo lô)}

Batch normalization chuẩn hóa từng giá trị trước kích hoạt (pre-activation) trên toàn bộ lô nhỏ (mini-batch) $B$, và được áp dụng ngay trước hoặc ngay sau hàm kích hoạt. Giải thuật Batch Normalization~\cite{ioffe2015batchnorm} gồm 4 bước áp dụng cho mỗi batch $B$:
\begin{enumerate}
    \item Tính trung bình của batch:
    \[ \boldsymbol\mu_B = \frac{1}{m_B} \sum_{i = 1}^{m_B} \mathbf{x}^{(i)}\]

    \item Tính phương sai của batch:
    \[ \boldsymbol\sigma^2_B = \frac{1}{m_B} \sum_{i = 1}^{m_B} \left(\mathbf{x}^{(i)} - \boldsymbol\mu_B \right)^2\]

    \item Chuẩn hóa
    \[ \hat{\mathbf{x}}^{(i)} = \frac{\mathbf{x}^{(i)} - \boldsymbol \mu_B}{\sqrt{\boldsymbol\sigma_B^2 + \varepsilon}}\]

    \item Scale và dịch chuyển
    \[ \mathbf{z}^{(i)} = \boldsymbol\gamma \otimes \hat{\mathbf{x}}^{(i)} + \boldsymbol\beta\]
\end{enumerate}
trong đó $\boldsymbol\mu_B$ và $\boldsymbol\sigma^2_B$ lần lượt là trung bình và phương sai của $\mathbf{x}_i$ trên toàn lô, $\varepsilon$ là một hằng số nhỏ để tránh chia cho 0, còn $\boldsymbol\gamma$ và $\boldsymbol\beta$ là các tham số có thể học được. Tại thời gian kiểm thử, giá trị trung bình động của $\boldsymbol\mu_B$ và $\boldsymbol \sigma_B^2$ sẽ được sử dụng.

Ioffe và Szegedy (2015)~\cite{ioffe2015batchnorm} nhận thấy rằng batch normalization đã cải thiện đáng kể hiểu quả của \textit{mọi} mạng sâu mà họ thử nghiệm, với mức cải thiện rất lớn trên bộ chuẩn ImageNet. Bài báo liệt kê những tác động sau:
\begin{itemize}
    \item BN giảm đáng kể hiện tượng gradient biến mất. Họ có thể sử dụng các hàm kích hoạt bão hòa như $\tanh$ và thậm chí sigmoid, những hàm trước đó được xem là gần như không thể sử dụng hiệu quả trong các mạng sâu.
    \item BN giảm đáng kể độ nhạy đối với khởi tạo trọng số. Phép chuẩn hóa tái căn giữa và tái điều chỉnh tỷ lệ của các tín hiệu ở mỗi lớp, nhờ đó một quy mô khởi tạo không phù hợp có thể được điều chỉnh.
    \item BN cho phép sử dụng tốc độ học lớn hơn nhiều, từ đó tăng đáng kể tốc độ học.
    \item BN giảm đáng kể số bước huấn luyện cần thiết. Trong bài báo, một mô hình phân loại hình ảnh tiên tiến nhất đã đạt cùng mức độ chính xác chỉ với số bước huấn luyện ít hơn 14 lần, và sau đó vượt xa mô hình gốc một khoảng cách lớn.
\end{itemize}

Tuy nhiên, các giá trị thống kê phụ thuộc vào lô dữ liệu, vì vậy các lô quá nhỏ sẽ tạo ra các ước lượng nhiễu, đồng thời cách tính toán lúc huấn luyện và lúc kiểm thử là khác nhau. Phương pháp này cũng khó áp dụng cho các mạng hồi quy (recurrent networks). Ngoài ra, Yang và cộng sự (2019)~\cite{yang2019} chứng minh rằng trong các mạng sâu có áp dụng BN nhưng không có residual connections, gradient có thể tăng theo mũ dựa trên chiều sâu mạng tại thời điểm khởi tạo, và điều này không thể khắc phục bằng cách thay đổi phương sai trọng số ban đầu hay hàm kích hoạt. Do đó, chuẩn hóa không thể thay thế cho việc thiết kế kiến trúc mạng tốt (Mục VII.3).

\subsubsection{Layer Normalization (Chuẩn hóa theo lớp)}

Layer Normalization, được đề xuất bởi Ba và cộng sự (2016)~\cite{ba2016layernorm}, tính toán các giá trị thống kê trên $n$ đơn vị của một lớp cho một mẫu dữ liệu đơn lẻ. Với một mẫu dữ liệu $\mathbf{x}$ có $n$ đặc trưng, ta tính
\[ \mu = \frac{1}{n}\sum_{j=1}^{n}x_j, \qquad \sigma^2=\frac{1}{n}\sum_{j=1}^{n}(x_j-\mu)^2, \qquad \mathbf{z}=\boldsymbol{\gamma}\otimes\frac{\mathbf{x}-\mu}{\sqrt{\sigma^2+\varepsilon}}+\boldsymbol{\beta}\]
Đây vẫn là quy trình bốn bước như batch normalization, nhưng thay vì tính thống kê trên mini-batch, layer normalization tính chúng dọc theo các đặc trưng của từng mẫu.

Vì không có yếu tố nào phụ thuộc vào lô dữ liệu (batch), layer normalization thực hiện cùng một phép tính tại thời điểm huấn luyện và kiểm thử, và có thể hoạt động tốt ngay cả với kích thước lô bằng 1. Vì vậy, layer normalization có thể áp dụng cho mạng hồi quy (recurrent network) bằng cách chuẩn hóa riêng biệt tại mỗi bước thời gian, và cách làm này giúp ổn định động lực của trạng thái ẩn. Đây là lựa chọn tiêu chuẩn trong các kiến trúc Transformer~\cite{vaswani2017}.

\subsection{Residual/skip connection và thiết kế kiến trúc}
\textit{Phụ trách: Tâm}

Giải thích đường tắt cho gradient, residual block và nguyên tắc thiết kế mạng sâu dễ tối ưu.

\subsection{Giải pháp dành cho RNN}
\textit{Phụ trách: Long}

LSTM, GRU, gradient clipping, truncated BPTT và lựa chọn kiến trúc thay thế khi phù hợp.

\subsection{Theo dõi gradient trong quá trình huấn luyện}
Đề xuất ghi log gradient norm, histogram trọng số/gradient và cảnh báo các lớp có gradient gần bằng $0$.

\begin{guidelinebox}{Thông điệp thực hành}
Không có một giải pháp duy nhất cho mọi mô hình. Nên kết hợp activation phù hợp, initialization đúng, normalization, residual connection và theo dõi gradient để chẩn đoán sớm.
\end{guidelinebox}

\section{Kết luận}
\textit{Người phụ trách: Trang, Trí}

\begin{itemize}
    \item Tóm tắt nguyên nhân cốt lõi: gradient là tích của nhiều đạo hàm/Jacobian nhỏ.
    \item Nhấn mạnh hậu quả: các lớp đầu hoặc các thời điểm xa không nhận đủ tín hiệu học.
    \item Tổng hợp dấu hiệu chẩn đoán và nhóm giải pháp phù hợp cho MLP/CNN/RNN.
    \item Nêu thông điệp cuối: thiết kế kiến trúc, khởi tạo và giám sát quá trình huấn luyện phải được xem xét đồng thời.
\end{itemize}

\clearpage
\renewcommand{\refname}{Tài liệu tham khảo}
\addcontentsline{toc}{section}{Tài liệu tham khảo}
\begin{thebibliography}{99}

\bibitem{fiedler2024}
Fiedler, K. (2024) `Olympische Spiele? Barenburg macht's ganz anders',
\textit{Kreiszeitung}, 12 August. Available at:
\url{https://www.kreiszeitung.de/lokales/diepholz/kirchdorf-ort120456/olympische-spiele-barenburg-macht-ganz-anders-spass-olympiade-93239133.html}
(Accessed: 30 September 2026).

\bibitem{hanin2018}
Hanin, B. (2018) `Which Neural Net Architectures Give Rise to Exploding and
Vanishing Gradients?', \textit{Advances in Neural Information Processing
Systems}, 31, pp.~580--589. Available at:
\url{https://proceedings.neurips.cc/paper_files/paper/2018/hash/13f9896df61279c928f19721878fac41-Abstract.html}
(Accessed: 30 September 2026).

\bibitem{he2016identity}
He, K., Zhang, X., Ren, S. and Sun, J. (2016) `Identity Mappings in Deep
Residual Networks', in Leibe, B., Matas, J., Sebe, N. and Welling, M. (eds.)
\textit{Computer Vision -- ECCV 2016}. Lecture Notes in Computer Science,
vol.~9908. Cham: Springer, pp.~630--645.
doi: \href{https://doi.org/10.1007/978-3-319-46493-0_38}{10.1007/978-3-319-46493-0\_38}.

\bibitem{krizhevsky2009}
Krizhevsky, A. (2009) \textit{Learning Multiple Layers of Features from Tiny
Images}. Technical report. Toronto: University of Toronto. Available at:
\url{https://www.cs.toronto.edu/~kriz/learning-features-2009-TR.pdf}
(Accessed: 30 September 2026).

\bibitem{rumelhart1986}
Rumelhart, D.E., Hinton, G.E. and Williams, R.J. (1986) `Learning
representations by back-propagating errors', \textit{Nature}, 323,
pp.~533--536. doi:
\href{https://doi.org/10.1038/323533a0}{10.1038/323533a0}.

\bibitem{bengio1994}
Bengio, Y., Simard, P. and Frasconi, P. (1994) `Learning long-term
dependencies with gradient descent is difficult', \textit{IEEE Transactions
on Neural Networks}, 5(2), pp.~157--166. doi:
\href{https://doi.org/10.1109/72.279181}{10.1109/72.279181}.

\bibitem{glorot2010}
Glorot, X. and Bengio, Y. (2010) `Understanding the difficulty of training
deep feedforward neural networks', in Teh, Y.W. and Titterington, M. (eds.)
\textit{Proceedings of the Thirteenth International Conference on Artificial
Intelligence and Statistics}. Proceedings of Machine Learning Research,
vol.~9, pp.~249--256. Available at:
\url{https://proceedings.mlr.press/v9/glorot10a.html}
(Accessed: 30 September 2026).

\bibitem{he2015rectifiers}
He, K., Zhang, X., Ren, S. and Sun, J. (2015) `Delving Deep into Rectifiers:
Surpassing Human-Level Performance on ImageNet Classification', in
\textit{Proceedings of the IEEE International Conference on Computer Vision},
pp.~1026--1034. doi:
\href{https://doi.org/10.1109/ICCV.2015.123}{10.1109/ICCV.2015.123}.

\bibitem{ioffe2015batchnorm}
Ioffe, S. and Szegedy, C. (2015) `Batch Normalization: Accelerating Deep
Network Training by Reducing Internal Covariate Shift', in Bach, F. and Blei,
D. (eds.) \textit{Proceedings of the 32nd International Conference on Machine
Learning}. Proceedings of Machine Learning Research, vol.~37, pp.~448--456.
Available at: \url{https://proceedings.mlr.press/v37/ioffe15.html}
(Accessed: 30 September 2026).

\bibitem{long-schoenholz2017}
Schoenholz, S.S., Gilmer, J., Ganguli, S. and Sohl-Dickstein, J. (2017) `Deep
Information Propagation', \textit{Proceedings of the 5th International
Conference on Learning Representations}. Available at:
\url{https://research.google/pubs/deep-information-propagation/}
(Accessed: 30 September 2026).

\bibitem{long-pascanu2013}
Pascanu, R., Mikolov, T. and Bengio, Y. (2013) `On the difficulty of training
recurrent neural networks', in Dasgupta, S. and McAllester, D. (eds.)
\textit{Proceedings of the 30th International Conference on Machine Learning}.
Proceedings of Machine Learning Research, vol.~28(3), pp.~1310--1318.
Available at: \url{https://proceedings.mlr.press/v28/pascanu13.html}
(Accessed: 30 September 2026).

\bibitem{long-goodfellow2016}
Goodfellow, I., Bengio, Y. and Courville, A. (2016) \textit{Deep Learning}.
Cambridge, MA: MIT Press. Available at:
\url{https://www.deeplearningbook.org/}
(Accessed: 30 September 2026).

\bibitem{borkar2025exploding}
V.~S. Borkar,
\newblock ``Exploding and vanishing gradients in deep neural networks: The effect of residual connections,''
\newblock \emph{arXiv preprint}, 2025.

\bibitem{ceni2022random}
A.~Ceni,
\newblock ``Random orthogonal additive filters: a solution to the vanishing/exploding gradient of deep neural networks,''
\newblock \emph{IEEE Transactions on Neural Networks and Learning Systems}, 2022.

\bibitem{poole2016}
Poole, B., Lahiri, S., Raghu, M., Sohl-Dickstein, J. and Ganguli, S. (2016)
`Exponential expressivity in deep neural networks through transient chaos',
in Lee, D., Sugiyama, M., Luxburg, U., Guyon, I. and Garnett, R. (eds.)
\textit{Advances in Neural Information Processing Systems}, 29.

\bibitem{yang2019}
Yang, G., Pennington, J., Rao, V., Sohl-Dickstein, J. and Schoenholz, S.S.
(2019) `A Mean Field Theory of Batch Normalization', \textit{Proceedings of
the 7th International Conference on Learning Representations}. Available at:
\url{https://openreview.net/forum?id=SyMDXnCcF7}
(Accessed: 6 October 2026).

\bibitem{ba2016layernorm}
Ba, J.L., Kiros, J.R. and Hinton, G.E. (2016) `Layer Normalization',
\textit{arXiv preprint} arXiv:1607.06450. Available at:
\url{https://arxiv.org/abs/1607.06450}
(Accessed: 6 October 2026).

\bibitem{vaswani2017}
Vaswani, A., Shazeer, N., Parmar, N., Uszkoreit, J., Jones, L., Gomez, A.N.,
Kaiser, L. and Polosukhin, I. (2017) `Attention Is All You Need', in
\textit{Advances in Neural Information Processing Systems}, 30,
pp.~5998--6008. Available at:
\url{https://arxiv.org/abs/1706.03762}
(Accessed: 6 October 2026).

\end{thebibliography}

\end{document}
